% Modernised reading — fonds Alexandre Grothendieck, shelfmark 76
% (pages 1-62, the whole folder).
%
% This is an INTERPRETATION, not a transcription. It re-reads the pages in
% today's notation and vocabulary, states the hypotheses the manuscript leaves
% implicit, and drops the critical apparatus so that the mathematics reads
% straight through. Everything the brackets and underlines carried moves into a
% footnote. It is held to being mathematically correct as it stands; where the
% manuscript is loose or mistaken, the true statement is what appears here and
% the footnote says what the page has. In any dispute of substance the
% transcription governs: whoever cites Grothendieck must cite batch-01.fr.tex
% (pp. 1-20), batch-02.fr.tex (pp. 21-40), batch-03.fr.tex (pp. 41-60) or
% batch-04.fr.tex (pp. 61-62).
%
% Pass: Opus 5.5 (claude-opus-5-5), 2026-10-10 — first pass, unchecked against the pages by a human.
% The folder was read whole, its four batches together; this is one document
% for the shelfmark. It was made on Opus 5.5 and has not been measured against
% the reference reading of folder 115 (made on Opus 5). The four
% transcriptions were themselves made in parallel, unchecked by a human, and
% pages 15-20 were checked against the facsimile only at whole-sheet
% resolution. The literature named in the footnotes is cited from memory and
% was not checked against the sources; no priority is inferred from it.
%
% Scope. Pages 1, 10, 14, 34, 44, 46, 48 are blank; page 43 is a cover
% (« 2-cartes cellulaires » struck, « cube combinatoire ») whose words head
% the cube section. Every other page is read here.
%
% Spine. One question, asked three ways, then a model case:
%   pp. 2-13    first draft, « Cartes coniques, cartes cellulaires, et ens. à
%               opérateurs »: a cell complex as a graded poset F_* with the
%               diamond axioms; its maximal flags as a set acted on by the
%               cartographic group; filtered spaces; the « n-variété
%               combinatoire » and the « n-complexe combinatoire » (the
%               cascade of p. 38); n-cartes coniques; flag lattices,
%               multicones, paracartes. Page 5 is cancelled.
%   pp. 15-33   « géométries de drapeaux »: posets whose lower ideals are
%               simplices (simplicial posets); dimension functions; total
%               orders on simplices = semi-simplicial sets with condition (1);
%               when the orders come from one order on the vertices; Ex 1
%               (order complexes), Ex 2 (G-sets, flags of type H, the Lemme);
%               presheaves on finite subsets of N.
%   pp. 35-42   second draft, « Cartes cellulaires n-dim et groupes
%               cartographiques »: n-admissible filtrations, categories F_n and
%               F_n°, the functors E_n, E_n°, Th. 1-7 (stated, not proved),
%               cascades, the multicone psi_n.
%   pp. 45, 47  asides: small sets as finite geometries and orders of
%               Chevalley groups; four « Sujets DEA ».
%   pp. 49-62   « cube combinatoire »: the cube of a set with free involution,
%               its face poset, flags, antipody, orientation, regular
%               polytopes, conjugacy classes of Aut(C_n), the 3-cube and its
%               four diagonals, the « cube gauche » (hemicube) and its
%               automorphisms.
%
% Notation fixed for the whole document. G_n is the n-dimensional
% cartographic group <sigma_0, ..., sigma_n | sigma_i^2, (sigma_i sigma_j)^2
% for |i-j| >= 2> (n+1 generators), as on pp. 7 and 36; page 3's G_n (n
% generators) is G_{n-1} here, page 13's G_i is G_{i-1}. sigma_i changes the
% i-dimensional member of a flag. « Drapeau » is any chain; a complete flag is
% « drapeau complet » (his « repère » on pp. 54-56). P*(I) the non-empty finite
% subsets of I. For a 1-simplex z of a semi-simplicial set, s(z) and t(z) are
% its first and last vertex; on a 2-simplex u, d_0 u, d_1 u, d_2 u are the faces
% opposite vertices 0, 1, 2 (his delta_i on pp. 22-24 are not the standard
% faces on D_1). On pp. 39-40 the page uses psi_n for two functors; here
% rho_n : G_n-sets -> (Casc)_n (p. 39, and p. 40's rho_{n-1}) and psi_n :
% (Casc)_n -> F_n (p. 40). Partial sections of B over I are Sp(B).
%
% Substantive departures from the pages, each footnoted where it occurs:
% - p. 3: G_n is presented without sigma_alpha^2 = 1; it is not a Coxeter group
%   of type A (consecutive generators satisfy no relation).
% - p. 2: « condition des chaînes » with chains of length <= n is not enough
%   for the poset to come from (F_alpha, I_alpha); the correct condition is
%   that every covering raises the height by one (example given).
% - p. 6: the « il faut et il suffit » for |F_*| to be a manifold is false as
%   written (suspension of a torus); the upper links must be spheres too.
% - p. 7: full faithfulness only on the isotopy category.
% - p. 13, p. 16: functors written covariant are contravariant (restriction).
% - p. 23: the correspondence is with (1) + transitivity (2°), not (2 bis).
% - p. 24: delta_0 for delta_1. p. 25: (a)-(d) characterise more than Ex 1
%   (boundary of a triangle); (a') + (d) does not give (a) (example).
% - p. 35: « n-admissible » for (n-1)-admissible. p. 37: sigma_1 ... sigma_{n-2}
%   for sigma_0 ... sigma_{n-2}.
% - p. 38, Th. 1: the essential image as stated is right only for n <= 2; for
%   n >= 3 cell boundaries and vertex links must be spheres (p. 42 says so).
% - p. 45: the formula for P_{G/B} is inverted.
% - p. 50: Isom_proj = Isom_aff fails for n = 1. p. 53: the centre of Aut(C) is
%   {1, sigma} for every n >= 1. p. 54: sigma_i, 0 <= i <= n-1 (not 1..n).
% - p. 55-56: « homogène » read with d_i counted downward, so that card Rep =
%   d_0 ... d_{n-1}.
% - p. 60: the injection is into Aut(Phi_0(C)/sigma).
% - p. 61: the intersection formula read p(F ∩ F') ∪ p(F ∩ sigma F').
% - p. 62: p(F) ⊂ p(F') iff F ⊂ F' or F ⊂ sigma F'. The « ? » on the composite
%   Aut(C)/sigma -> Aut_ord(Phi(C~)) is answered (yes, for all n >= 3) by this
%   reading, which says so.
%
% Announced and never established in the folder: Th. 1-4 and 7 (stated, no
% proof), Th. 5-6 (named only); the essential image of p. 4 and p. 7; whether
% the axioms imply bijectivity of A/G^(delta) -> A_delta (p. 3); the question
% of p. 19 (when the |F| form a cell decomposition); programme a), b) of p.
% 29; the case n even of p. 62.
\documentclass[11pt,a4paper]{article}
\input{../preamble/grothendieck}

\folder{76}
\pages{1}{62}
\dating{[vers 1977]}
\watermark{Édition de démonstration\\interprétation personnelle de l'œuvre}
\foldertitle{n-cartes cellulaires : notes manuscrites (s.d.) — lecture modernisée du dossier entier}

\begin{document}

\section*{Résumé}

\begin{resume}
Un cube, un ballon de football, un pavage de la sphère par des triangles :
ce sont des espaces découpés en morceaux simples --- des points, des arêtes,
des faces, des cellules de dimension supérieure --- recollés le long de leurs
bords. Ces feuillets cherchent à dire, de façon purement combinatoire, ce
qu'est un tel découpage, et à quel point la combinatoire suffit à retrouver
l'espace.

L'idée centrale est celle des \emph{drapeaux}. Dans un cube, un drapeau
complet est un choix d'un sommet, d'une arête qui le contient, d'une face qui
contient l'arête. Il y en a quarante-huit. Sur l'ensemble des drapeaux
agissent des opérations très simples : changer le sommet en gardant l'arête et
la face (il n'y a qu'une autre possibilité), changer l'arête en gardant le
sommet et la face, changer la face en gardant le sommet et l'arête. Chacune
est une involution, et deux opérations qui ne touchent pas des dimensions
voisines commutent. Le groupe ainsi défini, que Grothendieck appelle
\emph{groupe cartographique}, ne dépend que de la dimension ; un découpage
devient un ensemble sur lequel ce groupe agit, et l'on peut retrouver les
sommets, les arêtes et les faces comme les orbites de certains sous-groupes.
C'est l'analogue, en dimension quelconque, de ce que les dossiers 88 et 81
font pour les cartes tracées sur une surface.

Le dossier pose trois fois la même question, chaque fois plus nettement.
Une première rédaction donne des axiomes à l'ensemble ordonné des cellules ---
entre deux cellules dont les dimensions diffèrent de deux, il y en a toujours
exactement deux --- et montre comment les drapeaux le reconstituent. Une
longue étude intermédiaire porte sur les ensembles ordonnés dont chaque
élément « est » un simplexe : comment les décrire par des foncteurs, comment
les ordonner, comment passer d'une triangulation à une décomposition en
cellules. La seconde rédaction énonce un programme de sept théorèmes : les
variétés munies d'une bonne décomposition cellulaire, à déformation près,
sont exactement certains ensembles à opérateurs sous le groupe cartographique ;
une construction par cônes emboîtés, le « multicône », va dans l'autre sens.
Aucun de ces théorèmes n'est démontré, et Grothendieck voit lui-même où est
la difficulté : savoir si un objet combinatoire décrit une sphère, ce qui
touche à l'hypothèse de Poincaré.

Le dossier se termine sur un cas modèle traité à fond : le cube. Un cube de
dimension $n$ est décrit par un ensemble de $2n$ éléments --- ses faces de
codimension un --- muni de l'involution qui échange deux faces opposées. Tout
s'en déduit : les faces, les drapeaux, le groupe des symétries, qui agit
simplement transitivement sur les drapeaux, l'orientation, les classes de
conjugaison des symétries (dix pour le cube ordinaire), et le fait que les
rotations du cube ordinaire sont exactement les permutations de ses quatre
grandes diagonales. En recollant chaque face du bord avec la face opposée, on
obtient un découpage de l'espace projectif réel, et la question de ses
symétries ferme le dossier.

On voit dans ces pages une manière de procéder : plutôt que de choisir une
orientation, un ordre, un point de départ, garder l'ensemble de tous les
choix et faire agir sur lui le groupe qui les échange. Un cube « est » alors
un torseur sous son groupe, et une variété découpée « est » un ensemble à
opérateurs. Les noms modernes à chercher sont ceux de polytope abstrait, de
carte généralisée, d'ensemble ordonné simplicial et d'ensemble
semi-simplicial, de géométrie de classes à la Tits, et de groupe
hyperoctaédral.

\keywords{abstract polytope, diamond condition, flag graph, cartographic group, generalized map, order complex, simplicial poset, semi-simplicial set, category of elements, coset geometry, CW complex, Alexander trick, sphere recognition, hyperoctahedral group, signed cycle type, hemicube, Bruhat decomposition, exceptional isomorphisms of small groups}
\end{resume}

\pagerange{2}{62}
\section*{Le fil du dossier, et les conventions}

\subsection*{Les stations}

Le dossier n'est pas paginé de sa main. L'inventaire le date « vers 1977 » et
le range dans le groupe « Géométrie et topologie combinatoire », avec les
dossiers 81, 88 et 89, qui ont déjà leur lecture modernisée et qu'on cite là
où ils rencontrent celui-ci.

\begin{enumerate}
  \item \emph{Pages 2 à 13}, première rédaction, sous le titre « Cartes
  coniques, cartes cellulaires, et ensembles à opérateurs » : l'ensemble
  ordonné gradué $F_*$ des cellules et ses axiomes (pages 2 à 4) ; une page
  barrée de suites exactes (page 5) ; réalisation géométrique, variétés,
  « $n$-variété combinatoire » et « $n$-complexe combinatoire » (pages 6 à 8) ;
  « $n$-cartes coniques » (page 9) ; treillis de drapeaux, multicônes et
  « paracartes » (pages 11 à 13).
  \item \emph{Pages 15 à 33}, les « géométries de drapeaux » : ensembles
  ordonnés dont chaque idéal principal est un simplexe, fonctions dimension,
  ensembles semi-simpliciaux et ordres sur les sommets, deux exemples, et les
  préfaisceaux sur les parties finies de $\mathbf{N}$.
  \item \emph{Pages 35 à 42}, seconde rédaction, sous le titre « Cartes
  cellulaires $n$-dim et groupes cartographiques » : filtrations admissibles,
  les foncteurs $E_n$, $E_n^\circ$, et un programme de théorèmes numérotés 1 à
  7, énoncés et non démontrés.
  \item \emph{Pages 45 et 47}, deux apartés : les petits ensembles comme
  géométries finies et l'ordre des groupes de Chevalley ; quatre « Sujets
  DEA ».
  \item \emph{Pages 49 à 62}, le « cube combinatoire », sous une chemise
  (page 43) où « 2-cartes cellulaires » est biffé.
\end{enumerate}

Les deux rédactions des stations 1 et 3 ne se confondent pas. La première
part d'un ensemble ordonné et lui impose des axiomes ; la seconde part d'un
espace filtré et construit les ensembles à opérateurs par récurrence sur la
dimension. Elles se rejoignent sur deux objets, qu'on signale à chaque fois :
la tour d'ensembles à opérateurs des pages 7 et 8 est la « cascade » de la
page 38, et les quotients de la page 8 sont le théorème 3 de la page 39. On
garde les deux rédactions distinctes.

\subsection*{Conventions, valables pour tout le dossier}

\begin{itemize}
  \item \emph{Le groupe cartographique.} $G_n$ est le groupe engendré par
  $n+1$ éléments $\sigma_0, \ldots, \sigma_n$ soumis aux relations
  $\sigma_i^2 = 1$ et $(\sigma_i\sigma_j)^2 = 1$ pour $|i - j| \geqslant 2$,
  et à aucune autre ; $G_{-1} = \{1\}$. C'est la convention des pages 7 et 36.
  La page 3 note $G_n$ le groupe à $n$ générateurs $\sigma_0, \ldots,
  \sigma_{n-1}$, et la page 13 $G_i$ celui qu'engendrent $\sigma_0, \ldots,
  \sigma_{i-1}$ : ce sont ici $G_{n-1}$ et $G_{i-1}$.\footnote{Le nom
  « groupe cartographique » est le sien (page 35) ; c'est aussi celui de
  l'\emph{Esquisse d'un programme} (1984) pour $n = 2$, et le groupe
  $\widehat{\Gamma}$ de la page 10 du dossier 88 est $G_2$. C'est un groupe de
  Coxeter, de matrice $m_{ij} = 2$ pour $|i - j| \geqslant 2$ et $m_{i,i+1} =
  \infty$.}
  \item \emph{Drapeaux.} Un drapeau est une chaîne non vide d'un ensemble
  ordonné ; un \emph{drapeau complet} est une chaîne maximale. Les pages 54 à
  56 appellent « repères » les drapeaux complets d'un polyèdre, et la page 51
  les note $\mathrm{Drap}(C)$ ; on garde les deux notations, qui désignent la
  même chose. $\sigma_i$ change toujours le membre de dimension $i$ d'un
  drapeau complet en gardant les autres.
  \item \emph{Ensembles de parties.} $\mathfrak{P}^*(I)$ est l'ensemble des
  parties finies non vides de $I$, ordonné par inclusion.
  \item \emph{Semi-simplicial.} $\Delta_{\mathrm{inj}}$ est la catégorie des
  ensembles finis non vides totalement ordonnés et des applications
  strictement croissantes, $\Delta_n = [0, n]$ ; un \emph{ensemble
  semi-simplicial} est un préfaisceau sur $\Delta_{\mathrm{inj}}$. C'est ce que
  les pages appellent « ensemble quasi-simplicial » et notent
  « $\frac{1}{2}$simpl ».\footnote{On disait aussi « $\Delta$-ensemble »
  (Rourke et Sanderson, 1971). La page 20 note $\Delta$ cette catégorie ; on
  ajoute l'indice pour ne pas la confondre avec la catégorie simpliciale
  usuelle, qui a aussi les applications croissantes non injectives.} Pour un
  $1$-simplexe $z$, $s(z)$ et $t(z)$ sont son premier et son dernier sommet ;
  pour un $2$-simplexe $u$, $d_0u$, $d_1u$, $d_2u$ sont les arêtes opposées aux
  sommets $0$, $1$, $2$.
  \item \emph{Les deux foncteurs des pages 39 et 40.} La page 39 appelle
  $\psi_n$ un foncteur des $G_n$-ensembles vers les cascades, la page 40 un
  foncteur des cascades vers les espaces filtrés. On appelle ici $\rho_n$ le
  premier --- c'est la lettre que la page 40 elle-même emploie pour lui ---
  et $\psi_n$ le second.
  \item \emph{Le cube.} $(\mathbb{B}, \sigma)$ est un ensemble fini muni d'une
  involution sans point fixe, $I = \mathbb{B}/\sigma$, $n = \mathrm{card}\, I$,
  et $\mathrm{Sp}(\mathbb{B})$ l'ensemble des sections partielles de
  $\mathbb{B} \to I$, c'est-à-dire des parties $J \subset \mathbb{B}$ telles
  que $J \cap \sigma J = \emptyset$.
\end{itemize}

\emph{Ce que le dossier annonce et n'établit pas.} Les théorèmes 1 à 4 et 7
sont énoncés sans démonstration, les théorèmes 5 et 6 seulement nommés ;
« l'image essentielle » est à déterminer aux pages 4 et 7, et le reste ; la
question de la page 19 et le programme de la page 29 restent ouverts ; le
« cas $n$ pair » de la page 62 est « à examiner ».

\pagerange{2}{4}
\section*{Complexes cellulaires et ensembles à opérateurs (pages 2 à 4)}

\emph{Les données.} Des ensembles $F_0, \ldots, F_n$ --- les cellules de
dimension $0, \ldots, n$ --- et des relations d'incidence $I_\alpha \subset
F_\alpha \times F_{\alpha+1}$. Sur $F_* = \coprod_\alpha F_\alpha$ on définit
$x \preceq y$ si $x = y$, ou si $x \in F_\alpha$, $y \in F_\beta$, $\alpha <
\beta$, et s'il existe une chaîne $x = z_\alpha, z_{\alpha+1}, \ldots,
z_\beta = y$ d'éléments consécutivement incidents. Les axiomes, dans la
numérotation finale de la page 4 :\footnote{La numérotation est visiblement
remaniée sur la page : un « Axiome 3 » est surchargé en « 2 », « Axiome 4 »
est écrit sur un mot lourdement biffé, et la liste finale « $1, 1', 2, 2', 3,
4$ » est celle de la page 4. Le titre de la page 2 était « Complexes
cellulaires, et ens.\ à opérateurs » ; « Complexes » y est deux fois corrigé
en « Cartes », et « cellulaires » biffé au profit de « coniques ». La page 6
cite aussi un axiome « $3'$ » qui n'est défini nulle part.}

\begin{enumerate}
  \item[1.] $I_\alpha \to F_{\alpha+1}$ est surjectif : toute cellule a une
  cellule de dimension immédiatement inférieure dans son bord ;
  \item[1'.] (facultatif) $I_\alpha \to F_\alpha$ est surjectif ;
  \item[2.] pour $2 \leqslant \alpha \leqslant n$, $x_\alpha \in F_\alpha$ et
  $x_{\alpha-2} \preceq x_\alpha$ de dimension $\alpha - 2$, il y a exactement
  deux $x_{\alpha-1}$ entre eux ;
  \item[2'.] (facultatif) toute $x_{n-1}$ est incidente à exactement deux
  $x_n$ ;
  \item[3.] toute arête $x_1$ a exactement deux sommets ;
  \item[4.] voir plus bas.
\end{enumerate}

Les axiomes 2 et 3 forment ce qu'on appelle aujourd'hui la \emph{condition du
diamant} : tout intervalle de longueur deux, en ajoutant un plus petit élément
formel sous $F_0$, a exactement quatre éléments. L'axiome $2'$ est la même
condition en haut, et exprime l'absence de bord.\footnote{Ce sont, à la
connexité forte près (notre axiome 4), les axiomes des \emph{polytopes
abstraits} --- Danzer et Schulte (1982), puis McMullen et Schulte,
\emph{Abstract Regular Polytopes} (2002) ---, qui ajoutent aussi un plus
grand élément. Les feuillets ne citent rien de tel, et l'on ne tire de la
comparaison aucune conclusion d'antériorité : la datation de l'inventaire est
approximative.}

\emph{Quels ensembles ordonnés ?} Le système $(F_\alpha, I_\alpha)$ se
retrouve à partir de $(F_*, \preceq)$ : $F_0$ est l'ensemble des éléments
minimaux, $F_1$ celui des éléments minimaux de $F_* \smallsetminus F_0$, et
ainsi de suite, et $I_\alpha$ est la restriction de l'ordre. Les ensembles
ordonnés ainsi obtenus sont exactement ceux dont les chaînes ont au plus
$n+1$ éléments et dans lesquels \emph{tout recouvrement $x \lessdot y$ fait
monter la hauteur d'exactement un} --- la hauteur d'un élément étant le
nombre d'éléments, moins un, de la plus longue chaîne qui finit en lui ; de
façon équivalente, toutes les chaînes maximales de chaque idéal principal
ont la même longueur.\footnote{La page dit « exactement ceux satisfaisant la
condition des chaînes, et où les chaînes sont de longueur $\leqslant n$ », et
écrit « Vérifier » en marge. La condition de Jordan--Dedekind (dans chaque
intervalle, les chaînes maximales ont même longueur) ne suffit pas :
l'ensemble $\{a, b, d, c\}$ avec $a < c$ et $b < d < c$ la satisfait, mais
$c$ recouvre $a$ en montant de deux, et la reconstruction par niveaux
oublierait $a \preceq c$. Il faut la condition sur les idéaux principaux
qu'on énonce. Le même raisonnement montre que l'axiome 1 est alors
automatique, comme le dit le crochet de la page : « Cela montre pourquoi
$I_\alpha \to F_{\alpha+1}$ surjectif ».}

\emph{Drapeaux et opérateurs.} Un drapeau est une chaîne non vide de $F_*$ ;
son \emph{type} est l'ensemble $\delta \subset [0, n]$ des dimensions de ses
membres.\footnote{La page écrit d'abord « type $\in \mathfrak{P}[0,n]$ »,
puis entre parenthèses « une partie non vide de $[1,n]$ » ; le $0$ de
$[0,n]$ est repassé à l'encre noire. Le type d'un drapeau est une partie non
vide de $[0, n]$.} Soit $\mathcal{A}$ l'ensemble des drapeaux complets $d =
(x_0, x_1, \ldots, x_n)$, $x_i \in F_i$. Par les axiomes 2 et 3, pour $0
\leqslant \alpha \leqslant n-1$, il existe un unique $x'_\alpha \neq
x_\alpha$ tel que $\sigma_\alpha(d) = (x_0, \ldots, x'_\alpha, \ldots, x_n)$
soit encore un drapeau ; sous l'axiome $2'$, de même pour $\alpha = n$. Ce
sont des involutions sans point fixe, et
\[
  \sigma_\alpha\sigma_\beta = \sigma_\beta\sigma_\alpha \qquad \text{si }
  |\alpha - \beta| \geqslant 2 ,
\]
puisqu'elles modifient alors des membres non voisins du drapeau. Ainsi
$\mathcal{A}$ est un ensemble à opérateurs sous $G_{n-1}$, et sous $G_n$ si
l'axiome $2'$ est satisfait.\footnote{La page définit « $G_n$ le groupe
engendré par les gén.\ $\sigma_\alpha$ ($0 \leqslant \alpha \leqslant n-1$)
avec les relations précédentes », c'est-à-dire avec les seules relations de
commutation. Il faut y joindre $\sigma_\alpha^2 = 1$ pour que l'action se
factorise, et c'est le groupe des pages 7 et 36. On l'appelle $G_{n-1}$,
conformément à la convention fixée plus haut. Ce n'est pas un groupe de
Coxeter de type $A$ : deux générateurs consécutifs n'y satisfont aucune
relation, alors que dans le groupe symétrique $(\sigma_\alpha
\sigma_{\alpha+1})^3 = 1$.}

\emph{L'axiome 4.} Pour $\delta \subset [0, n]$ non vide, soit
$G^{(\delta)}$ le sous-groupe engendré par les $\sigma_i$, $i \notin
\delta$ : il agit sur $\mathcal{A}$ sans toucher aux membres de dimension
dans $\delta$, d'où une application
\[
  \mathcal{A}/G^{(\delta)} \longrightarrow \mathcal{A}_\delta
  \qquad (\text{drapeaux de type } \delta),
\]
surjective par l'axiome 1 et sa conséquence que tout drapeau se complète.
L'axiome 4 demande qu'elle soit bijective pour $\delta = \{\alpha\}$ :
\[
  \mathcal{A}/G^{(\alpha)} \xrightarrow{\ \sim\ } F_\alpha ,
\]
c'est-à-dire que deux drapeaux complets qui ont le même membre de dimension
$\alpha$ se déduisent l'un de l'autre par des $\sigma_\beta$, $\beta \neq
\alpha$. La page demande si la bijectivité pour tout $\delta$ en est
conséquence.\footnote{« conséquence \ldots\ de l'axiome 4 et des
précédents ? », avec un point d'interrogation. Elle l'est si l'axiome 4 vaut
pour chaque \emph{section} $[x, y]$ de $F_*$, et non seulement pour $F_*$ :
c'est la \emph{connexité forte par drapeaux} des polytopes abstraits. La page
note d'ailleurs que les $F(x) = \{y \prec x\}$ satisfont les mêmes axiomes,
avec $n$ remplacé par la dimension de $x$ moins un, ce qui donne les sections
inférieures ; les sections supérieures ne sont pas évoquées. Remarque de
notre part.} Comme $G^{(\delta)} \subset G^{(\delta')}$ pour $\delta' \subset
\delta$, les inclusions entre ces sous-groupes redonnent alors les relations
entre drapeaux de types divers, donc $F_*$ tout entier :

\textbf{Proposition} (page 4). Le foncteur $F_* \mapsto \mathcal{A}$, des
complexes satisfaisant les axiomes 1 à 4 vers les ensembles à opérateurs sous
le groupe cartographique, est pleinement fidèle ; son image essentielle reste
à déterminer.\footnote{La page porte « (confirmé ?) » à côté de la liste des
axiomes. Les morphismes de complexes ne sont pas définis ; l'énoncé est
juste si l'on prend pour morphismes les applications croissantes qui
respectent la dimension et induisent des bijections sur les drapeaux
complets de chaque section --- les revêtements --- puisque $F_\alpha$ se lit
alors comme $\mathcal{A}/G^{(\alpha)}$ de façon fonctorielle. C'est notre
précision. Les ensembles à opérateurs obtenus sont des « cartes
généralisées » au sens de la note de la page 7.}

\emph{Réalisation géométrique.} Celle de $F_*$ --- sous le seul axiome 1 ---
est la réalisation de la triangulation dont les simplexes sont les parties
finies totalement ordonnées de $F_*$ : le \emph{complexe des chaînes} (ou
complexe d'ordre) de $F_*$.

\pagerange{5}{5}
\section*{Une page barrée (page 5)}

La page 5 est barrée d'un long trait et ne porte aucun texte : des suites
exactes empilées de groupes d'automorphismes d'un espace $X$ à bord. On y
lit
\[
  1 \to \mathrm{Aut}(X, \mathrm{id}_{\partial X}) \to \mathrm{Aut}(X) \to
  \mathrm{Aut}(\partial X) \to 1 ,
\]
un revêtement $p : \widetilde{\partial X} \to \partial X$ et le groupe
$\widetilde{\mathrm{Aut}}(X)$ des couples $(u, \tilde v)$ où $\tilde v$
relève au revêtement la restriction de $u$ au bord, et un facteur
$\mathbf{Z}^I$, $I = \pi_0(\partial X)$, qui revient dans trois lignes
successives. Rien ne dit ce que sont les groupes notés $S\mathcal{A}$,
$S\mathcal{A}^\circ$, $ST$, ni si les « Aut » sont des groupes
d'homéomorphismes ou leurs classes d'isotopie ; on ne leur donne pas de
sens.\footnote{Le facteur $\mathbf{Z}^I$, un entier par composante du bord,
évoque les twists de Dehn le long des composantes du bord dans un groupe de
classes d'isotopie ; c'est une conjecture de lecture, que rien sur la page ne
confirme. La page porte aussi deux annotations, « $\simeq S\mathcal{A}(X)
\times \mathbf{Z}^I$ » et « $\simeq ST(X) \times \mathbf{Z}^I$ », sous les
groupes relevés.}

\pagerange{6}{8}
\section*{Réalisation, variétés, et la première « variété combinatoire » (pages 6 à 8)}

\emph{Les cellules dans la réalisation.} Chaque $x \in F_\alpha$ se réalise
dans $|F_*|$ comme la réunion des simplexes ouverts des chaînes formées
d'éléments $\preceq x$, et l'incidence devient l'inclusion. La page annonce
une équivalence de catégories « à dégager » entre ces objets topologiques et
les objets combinatoires.

\textbf{Proposition} (page 6). Si $|F_*|$ est une variété topologique de
dimension $n$, sans bord, alors $F_*$ satisfait les axiomes $1, 1', 2, 2', 3,
4$.\footnote{C'est vrai et se vérifie sur les links : une variété est une
pseudo-variété pure, donc toutes les chaînes maximales ont $n+1$ éléments ;
une chaîne à laquelle il manque exactement une dimension est une face de
codimension un du complexe des chaînes, contenue dans exactement deux
simplexes maximaux, ce qui donne les axiomes $2$, $2'$ et $3$ ; l'axiome 4
vient de ce que le link d'un sommet, une sphère homologique de dimension
$n-1$, est fortement connexe. Le « pointée » ajouté en interligne après
« variété top. » est d'une lecture douteuse ; on ne le reprend pas.}

La page affirme plus précisément que $|F_*|$ est une variété de dimension $n$
si et seulement si $F_*$ satisfait ces axiomes et que, pour tout $x$, la
réalisation de $F(x) = \{y \prec x\}$ est une sphère ; et, « modulo
l'hypothèse de Poincaré », qu'il suffirait de demander $|F(x)|$ simplement
connexe et de l'homologie d'une sphère. Ce n'est pas suffisant tel quel. La
condition sur les $|F(x)|$ fait de chaque cellule une boule, mais elle ne
dit rien de ce qui entoure une cellule : il faut aussi que les réalisations
des sections supérieures $\{y \succ x\}$, les links des cellules, soient des
sphères de la bonne dimension.\footnote{Contre-exemple : la suspension d'un
tore triangulé est un complexe simplicial de dimension 3 qui satisfait tous
les axiomes de la page 3, et dont toutes les cellules sont des simplexes,
donc des boules ; mais le link de chacun des deux points de suspension est un
tore, et l'espace n'est pas une variété. La page cite un axiome « $3'$ »
qu'elle ne définit pas ; si c'était l'énoncé dual sur les sections
supérieures, l'équivalence redeviendrait juste dans sa forme « sphère ».
Quant à la forme homologique : une sphère d'homologie simplement connexe qui
est une variété est une sphère, en toute dimension --- c'est l'hypothèse de
Poincaré généralisée, démontrée par Smale en dimension $\geqslant 5$ (1961),
Freedman en dimension 4 (1982), Perelman en dimension 3 (2003) ; vers 1977
seuls les cas $\geqslant 5$ l'étaient, et la page corrige justement
« $\alpha \geqslant 3$ » en « $\alpha \geqslant 5$ » dans son addition. Mais il
faut encore savoir que $|F(x)|$ est une variété, ce qui n'est pas
automatique. Le détail des surcharges de la page n'est pas sûr.}

\emph{Espaces filtrés.} Soit $X$ un espace --- compact, ajoute-t-il avec deux
points d'interrogation --- filtré par $X_0 \subset X_1 \subset \cdots \subset
X_n = X$, chaque différence $X_i \smallsetminus X_{i-1}$ étant une réunion
disjointe de boules ouvertes de dimension $i$. Les drapeaux relatifs à la
filtration, « à définir avec soin », forment un ensemble sur lequel agit
$G_n$ si $X$ est une variété de dimension $n$, les $\sigma_i$ opérant sans
point fixe.\footnote{La fin de la page 6 est serrée et surchargée ; « la
différence de $X_i$ moins $X_{i-1}$ » est suivie de mots illisibles, puis de
« de boules $B^i$ ». On lit, comme à la page 35, des boules ouvertes de
dimension $i$.}

\textbf{Proposition} (page 7). Le foncteur des variétés de dimension $n$
ainsi filtrées vers les $G_n$-ensembles finis où les $\sigma_i$ opèrent sans
point fixe est pleinement fidèle --- à condition de remplacer les morphismes
par leurs classes d'isotopie. Il reste à déterminer son image essentielle,
dont la page prévoit qu'elle est plus grande que celle des complexes de la
page 4.\footnote{La page dit « pl.\ fidèle » sans parler d'isotopie. Sans
passer aux classes d'isotopie le foncteur n'est pas fidèle : un homéomorphisme
isotope à l'identité à travers des homéomorphismes respectant la filtration
fixe tous les drapeaux. La seconde rédaction (page 38) passe en effet à la
« catégorie isotopique ». Pour l'image essentielle, voir la note au
théorème 1.}

\emph{La $n$-variété combinatoire.} Une \emph{$n$-variété combinatoire finie}
est un ensemble fini $D_n$ muni de $n+1$ involutions sans point fixe
$\sigma_0, \ldots, \sigma_n$, deux involutions non consécutives commutant :
c'est un $G_n$-ensemble fini où les générateurs opèrent sans point
fixe.\footnote{C'est exactement ce qu'on appelle aujourd'hui une
\emph{$n$-carte généralisée} sans bord (P.~Lienhardt, 1989--1994), ou une
carte combinatoire au sens de A.~Vince (1983) ; on peut aussi le voir comme un
graphe régulier de degré $n+1$ à arêtes colorées, voisin des
« cristallisations » de M.~Pezzana et de l'école de Modène (années 1970), qui
n'imposent pas les relations de commutation. Le $D_n$ de la page joue le rôle
de l'ensemble des drapeaux complets. Les pages ne citent aucun de ces
travaux.}

\emph{Le $n$-complexe combinatoire.} Un \emph{$n$-complexe combinatoire
fini} est une tour d'ensembles finis
\[
  D_n \longrightarrow D_{n-1} \longrightarrow \cdots \longrightarrow D_1
  \longrightarrow D_0 ,
\]
où chaque $D_i$, $1 \leqslant i \leqslant n$, est un $G_{i-1}$-ensemble dont
les générateurs opèrent sans point fixe, où chaque $D_i \to D_{i-1}$ est un
morphisme de $G_{i-2}$-ensembles, et où, pour $i \geqslant j + 2$, la
composée $D_i \to D_j$ est invariante par $\sigma_{j+1}, \ldots,
\sigma_{i-1}$. Pensez à $D_i$ comme à l'ensemble des drapeaux complets
$(x_0, \ldots, x_i)$ des cellules de dimension $i$, et à $D_i \to D_{i-1}$
comme à l'oubli de $x_i$. C'est la « cascade » que la seconde rédaction
définira à la page 38 ; on y revient.

\emph{Les bords successifs.} Si $D_n$ est une $n$-variété combinatoire, le
bord de la $i$-ème décomposition, $\partial(\mathrm{Déc}(X_i/X_{i-1}))$, est
la $(i-1)$-variété combinatoire
\[
  \delta_{i-1} = \bigl(D_n/\langle \sigma_n, \ldots, \sigma_{i+1} \rangle
  ;\ \sigma_0, \ldots, \sigma_{i-1}\bigr) ,
  \qquad \delta_{n-1} = (D_n ;\ \sigma_0, \ldots, \sigma_{n-1}) .
\]
Les $\sigma_j$, $j \leqslant i - 1$, passent au quotient parce qu'ils
commutent à $\sigma_{i+1}, \ldots, \sigma_n$. C'est la forme première du
théorème 3 de la page 39. La suite de la page 8, qui commençait à définir par
récurrence une catégorie $\mathcal{C}_n$ de « $n$-cartes » et un foncteur vers
les $n$-variétés combinatoires, est barrée.\footnote{Dans la partie barrée :
$\mathcal{C}_0$ = ensembles finis (« 0-sphères »), un objet de
$\mathcal{C}_n$ comme un triple $(X_{n-1}, S_{n-1}, \varphi)$ avec
$\varphi : i_{n-1}(S_{n-1}) \to X_{n-1}$ --- une cellule attachée par son bord
---, et le foncteur $i_n$ écrit tantôt dans un sens, tantôt dans l'autre.
C'est le schéma de la récurrence des pages 35 à 40.}

\pagerange{9}{9}
\section*{Cartes coniques (page 9)}

La page reprend le titre, corrigé : « Cartes coniques, cartes cellulaires, et
ensembles à opérateurs ». Une \emph{$n$-carte conique} finie est un espace
compact $X$ muni d'une filtration croissante par des fermés
$X_0 \subset \cdots \subset X_n = X$ telle que :
\begin{enumerate}
  \item[a)] $X_0$ est fini ;
  \item[b)] pour $1 \leqslant i \leqslant n$, soit $\overline{Z}_i =
  \mathrm{Déc}(X_i/X_{i-1})$ le \emph{découpage} de $X_i$ le long de
  $X_{i-1}$ et $\dot Z_i \subset \overline{Z}_i$ l'image inverse de $X_{i-1}$,
  de sorte que $\overline{Z}_i \to X_i$ induise un homéomorphisme
  $\overline{Z}_i \smallsetminus \dot Z_i \simeq X_i \smallsetminus X_{i-1}$ ;
  alors $\overline{Z}_i$ a un nombre fini de composantes connexes
  $\overline{Z}_i^\alpha$, chacune homéomorphe au cône sur $\dot Z_i^\alpha =
  \overline{Z}_i^\alpha \cap \dot Z_i$, non vide ;
  \item[c)] $\dot Z_i \to X_{i-1}$ est un morphisme fini ;
  \item[d)] la filtration est équisingulière.
\end{enumerate}
Le découpage n'est pas défini ; on le comprend comme l'espace des cellules
fermées, avec leurs applications d'attachement, ce qui est le sens que lui
donne la page 35.\footnote{La condition d) est récrite en surcharge à droite :
« soit $X_{i-1}^!$ l'image de $\dot Z_i$ dans $X_{i-1}$ ; alors $X_{i-2}^! =
X_{i-1}^! \cap X_{i-2}$ », lecture incertaine des indices. Dans c), un
complément entre crochets est illisible.}

La différence avec les cartes \emph{cellulaires} est dans b) : chaque
« cellule » est un cône sur son bord, et non une boule ; le bord n'a pas à
être une sphère. On retrouvera ces cônes dans le multicône de la page 40, et
c'est ce qui donne son titre à la rédaction. Pour $n = 0$ ce sont les espaces
finis discrets, pour $n = 1$ les graphes finis --- boucles et arêtes
multiples admises ---, dont la page dessine quelques exemples.

\pagerange{11}{13}
\section*{Treillis de drapeaux, multicônes, paracartes (pages 11 à 13)}

\emph{Treillis de drapeaux.} Pour une décomposition de dimension $\leqslant
n$, notons $D_I$ l'ensemble des drapeaux de type $I \subset [0, n]$ ; pour $J
\subset I$, l'oubli des membres de dimension hors de $J$ donne une
restriction $D_I \to D_J$. Les pages 11 et 13 dessinent ces treillis pour $n
= 1, 2, 3$ et marquent « épi » certaines flèches. La règle qui s'en dégage :
\[
  D_{[0,i]} \longrightarrow D_I \quad \text{est surjective si } i = \sup I ,
\]
tout drapeau se complétant vers le bas jusqu'au sommet, et plus généralement
$D_I \to D_J$ l'est si $J \subset I$ a le même plus grand élément. La flèche
$D_{02} \to D_1$ est barrée d'une croix sur le dessin : il n'y a pas de
restriction de $D_{\{0,2\}}$ vers $D_{\{1\}}$.\footnote{La page 11 écrit
aussi « $D_\sigma \to D_{\sigma'}$ épi si $\sigma' \subset \sigma$ et
\ldots, i.e.\ $\sigma = \sigma'' \cup \sigma'$, $\sigma'' < \sigma'$ » : $J$ doit
être un segment final de $I$, ce qui est la même règle. La formule $n_i =
\sum_{j<i} n_j + 1$ qui termine la page, à côté d'un dessin d'arcs hachurés,
reste sans contexte ; on ne l'interprète pas.} La page 12 donne le cas d'un
graphe : trois sortes de drapeaux, les sommets $F_0$, les arêtes $F_1$, et les
couples incidents $R_1$, avec $F_0 \leftarrow R_1 \to F_1$, la seconde flèche
surjective.

\emph{Le multicône.} Pour une catégorie $\mathcal{C}$, soit $\Phi(\mathcal{C})$
la catégorie des familles finies $(f_i : X_i \to X)_{i \in I}$ de morphismes
de même but. Un foncteur $\varphi : \mathcal{C} \to \mathrm{Top}$ en induit
un, $\Phi(\varphi) : \Phi(\mathcal{C}) \to \mathrm{Top}$,
\[
  (X_i \xrightarrow{f_i} X)_{i \in I} \longmapsto
  \Bigl(\coprod_i \mathrm{Cône}\, \varphi(X_i)\Bigr)
  \ \amalg_{\coprod_i \varphi(X_i)}\ \varphi(X) ,
\]
qui colle à $\varphi(X)$ un cône sur chaque $\varphi(X_i)$ le long de
$\varphi(f_i)$. Si les $\varphi(f)$ sont finis, les $\Phi(\varphi)(f)$ le
sont ; la page ajoute que la fidélité --- et la pleine fidélité, avec un
doute --- passe de $\varphi$ à $\Phi(\varphi)$ vers les espaces filtrés,
« mod isotopie », et qu'on peut itérer : $\Phi(\Phi(\mathcal{C})) \to
\mathrm{Top}$. Itéré $n$ fois, c'est exactement la construction des
$n$-cartes coniques de la page 9, et c'est le « multicône » de la page
40.\footnote{Le rapprochement avec la page 40 est le nôtre ; la page 12 ne
nomme pas la construction. Sur ce que deviennent les morphismes de
$\Phi(\mathcal{C})$, la page ne dit rien ; une flèche $(X_i \to X) \to (Y_j
\to Y)$ doit comporter une application d'ensembles d'indices et des
morphismes compatibles, et l'énoncé de fidélité dépend de ce choix.}

\emph{Paracartes} (page 13). Une \emph{paracarte} est un préfaisceau $D$ sur
$\mathfrak{P}^*(\mathbf{N})$ --- un ensemble $D(\sigma)$ de « drapeaux de
type $\sigma$ » pour chaque partie finie non vide $\sigma$ de $\mathbf{N}$,
avec des restrictions $D(\sigma) \to D(\sigma')$ pour $\sigma' \subset
\sigma$ --- dont certaines restrictions sont surjectives.\footnote{La page
écrit « foncteurs $\mathfrak{P}_f^*(\mathbf{N}) \xrightarrow{D} (\mathrm{Ens})$
tels que $\sigma \subset \sigma'$ [segments finis] $\Rightarrow D(\sigma) \to
D(\sigma')$ épi », la flèche allant de la plus petite partie vers la plus
grande. Un drapeau de type $\sigma$ se restreint en un drapeau de type $\sigma'
\subset \sigma$, et non l'inverse ; le foncteur est contravariant, comme il
l'est aux pages 15 et 30, et on rétablit le sens des flèches. Le premier
$\sigma$ de la formule est surchargé, peut-être sur un $\Delta$.} La page en
distingue trois sortes :
\begin{itemize}
  \item les \emph{stratifications équidimensionnelles} : $D(\sigma) \to
  D(\sigma')$ est surjective dès que $\sup \sigma = \sup \sigma'$,
  c'est-à-dire que $D([0, i]) \to D(\sigma)$ l'est si $i = \sup \sigma$ ---
  la règle des treillis ci-dessus ;
  \item les paracartes \emph{de dimension $n$} : les restrictions sont
  bijectives au-dessus de $\Delta_n$ et $D(\{i\}) = \emptyset$ pour $i > n$ ;
  \item les paracartes \emph{spéciales} : équidimensionnelles, et telles que
  pour tout $j < i$ la restriction $D(0, \ldots, i) \to D(0, \ldots,
  \widehat{j}, \ldots, i)$ soit de degré 2 --- la condition du diamant.
\end{itemize}
Dans une paracarte spéciale, $R_i = D(0, \ldots, i)$ porte donc $i$
involutions $\sigma_0, \ldots, \sigma_{i-1}$ : c'est un
$G_{i-1}$-ensemble.\footnote{La page écrit « groupe $G_i$ engendré par
$\sigma_0, \ldots, \sigma_{i-1}$ opérant sur $R_i$ » ; c'est notre $G_{i-1}$,
conformément à la convention. Il n'y a pas de $\sigma_i$ sur $R_i$ : changer
le membre de dimension $i$ demande qu'une cellule de dimension $i - 1$ soit
dans exactement deux cellules de dimension $i$, ce qu'une paracarte spéciale
n'exige pas. La page demande aussi si la restriction $D(\ldots, j-1, j, j+1,
\ldots) \to D(\ldots, j-1, j+1, \ldots)$, pour un type qui n'est pas un
segment initial, est toujours de degré 2 (« p.\ ex.\ $D_{123} \to D_{13}$ ») ;
la question reste ouverte.} La page veut des conditions pour que $D(0,
\ldots, i)$, avec l'action de $G_{i-1}$ et les applications $R_{i+1} \to R_i
\to \cdots$, « récupère » tous les $D(i_0, \ldots, i_p)$ : c'est la tour de la
page 7, et la question à laquelle répondra le théorème 3.

\pagerange{15}{19}
\section*{Géométries de drapeaux (pages 15 à 19)}

\emph{Définition.} Un ensemble ordonné $D$ est une \emph{géométrie de
drapeaux} si, pour tout $d \in D$, l'idéal $D(d) = \{d' \leqslant d\}$ est
isomorphe à $\mathfrak{P}^*(I_d)$ pour un ensemble fini non vide $I_d$. On
peut prendre pour $I_d$ l'ensemble des éléments minimaux de $D$ majorés par
$d$ ; on appelle $\Phi$ l'ensemble des éléments minimaux de $D$ --- les
« drapeaux de longueur 1 » --- et $\mathrm{card}\, I_d$ la longueur de
$d$.\footnote{On dit aujourd'hui \emph{ensemble ordonné simplicial}
(A.~Björner, 1984 ; Garsia et Stanton, 1984), en ajoutant d'ordinaire un plus
petit élément, la face vide : chaque idéal principal est un treillis
booléen. Un élément est un simplexe, mais deux simplexes peuvent avoir les
mêmes sommets. Les pages ne citent pas ces travaux, postérieurs à la
datation de l'inventaire.}

Pour $\sigma \in \mathfrak{P}^*(\Phi)$, soit $D(\sigma)$ l'ensemble des $d$
tels que $I_d = \sigma$. Si $\tau \subset \sigma$, chaque $d \in D(\sigma)$ a
un unique minorant dans $D(\tau)$, d'où une restriction $D(\sigma) \to
D(\tau)$. On obtient :

\textbf{Proposition} (page 15). Les géométries de drapeaux sont
équivalentes aux couples $(\Phi, D)$ formés d'un ensemble $\Phi$ et d'un
préfaisceau $D : \mathfrak{P}^*(\Phi)^{\mathrm{op}} \to \mathrm{Ens}$ tel que
$\mathrm{card}\, D(\{i\}) = 1$ pour tout $i$ ; on retrouve $D$ comme
$\coprod_\sigma D(\sigma)$, avec $d \leqslant d'$ si $d$ est une restriction
de $d'$.

Une géométrie de drapeaux est \emph{stricte} si tout $d$ est la borne
supérieure de ses sommets. Il faut et il suffit pour cela que $\mathrm{card}\,
D(\sigma) \leqslant 1$ pour tout $\sigma$, c'est-à-dire que $D$ soit un
sous-foncteur du foncteur final ; $D$ est alors un ensemble de parties finies
non vides de $\Phi$ stable par passage aux parties non vides, et :
\[
  \{\text{géométries de drapeaux strictes}\} \simeq
  \{\text{complexes simpliciaux abstraits}\} .
\]
\footnote{La page écrit « ensembles simpliciaux » (lecture incertaine du
premier mot) ; dans l'usage d'aujourd'hui un ensemble simplicial est tout
autre chose, et ce sont les complexes simpliciaux abstraits --- les
« multiplicités simpliciales » de la page 25 --- qui conviennent. Une note de
marge ajoute que $\Phi$ est aussi l'ensemble des sommets d'un ensemble
semi-simplicial, mais que $D$ n'est pas déterminé par celui-ci.}

\emph{Fonctions dimension} (pages 16 et 17). Partons maintenant d'un
préfaisceau quelconque $D' : \mathfrak{P}^*(I)^{\mathrm{op}} \to
\mathrm{Ens}$, sans condition sur les $D'(\{i\})$. L'ensemble $D =
\coprod_\sigma D'(\sigma)$, ordonné comme plus haut, est une géométrie de
drapeaux, d'ensemble de sommets $\Phi = \coprod_{i \in I} D'(\{i\})$, et la
projection $\delta : \Phi \to I$ est injective sur chaque $I_d$.
Réciproquement, d'une géométrie de drapeaux $D$ et d'une application
$\delta : \Phi \to I$ injective sur chaque $I_d$, on tire
\[
  D'(\sigma) = \coprod_{\substack{\sigma' \in \mathfrak{P}^*(\Phi) \\
  \delta|\sigma' : \sigma' \xrightarrow{\sim} \sigma}} D(\sigma') ,
\]
et l'on a une équivalence
\[
  \bigl(D, \ \delta : \Phi \to I\bigr) \quad \longleftrightarrow \quad
  \bigl(D' : \mathfrak{P}^*(I)^{\mathrm{op}} \to \mathrm{Ens}\bigr) .
\]
\footnote{La page 16 écrit $D' : \mathfrak{P}_f^*(I) \to (\mathrm{Ens})$ sans le
signe d'opposé qu'elle met partout ailleurs ; c'est le même préfaisceau. La
page 17 ajoute que $I_0 = \delta(\Phi)$, l'ensemble des $i$ tels que
$D'(\{i\}) \neq \emptyset$, porte une structure d'ensemble semi-simplicial
« \ldots\ $\sigma$ tel que $D(\sigma) \neq \emptyset$ », dont la fin est
illisible.} Le cas qui intéresse la suite est $I = \mathbf{N}$, ou $I = [0,
n]$ : $\delta$ est alors une \emph{fonction dimension}, et $\sup \delta(\Phi)$,
si elle est finie, la dimension de $D$. Ainsi, une géométrie de drapeaux munie
d'une fonction dimension n'est rien d'autre qu'une paracarte au sens de la
page 13.

\emph{Ordres sur les simplexes} (pages 17 et 18). Supposons donné sur chaque
simplexe $I_d$ un ordre total, ces ordres s'induisant mutuellement --- c'est
le cas s'ils viennent d'une fonction dimension à valeurs dans un ensemble
totalement ordonné. Posons $D[m] = \{d \mid \mathrm{card}\, I_d = m + 1\}$.
Alors $m \mapsto D[m]$ est un ensemble semi-simplicial : une application
strictement croissante $u : \Delta_k \to \Delta_m$ choisit, dans le simplexe
ordonné $I_d$, une face, et $u^*(d)$ est la restriction de $d$ à cette
face.\footnote{La page écrit « un foncteur sur la catégorie $\Delta$ » ; le
« contra » est restitué dans la transcription. La phrase entre crochets sur la
façon dont un ordre de $I$ induit ces ordres totaux est en partie
illisible.} Cet ensemble semi-simplicial satisfait :

\begin{enumerate}
  \item[(1)] pour $k < m$ et $d \in D[m]$, l'application
  $\mathrm{Hom}(\Delta_k, \Delta_m) \to D[k]$, $u \mapsto u^*(d)$, est
  injective : les faces d'un simplexe sont deux à deux distinctes.
\end{enumerate}

La condition (1) équivaut à ce que la catégorie des éléments
$\Delta_{\mathrm{inj}}/D$ soit un ensemble ordonné, et cet ensemble ordonné
est $D$.\footnote{La page écrit « la catégorie $\widehat{\Delta}/D[\,]$ est
ordonnée (i.e.\ équivalente \ldots) ». La catégorie des éléments est
$\Delta/D$, sous-catégorie pleine de $\widehat{\Delta}/D$ par Yoneda ; c'est
d'elle qu'il s'agit. La page 20 ajoute qu'« ordonnée » et « préordonnée »
reviennent ici au même : c'est exact, parce que les seuls endomorphismes de
$\Delta_{\mathrm{inj}}$ sont les identités, de sorte que deux objets
isomorphes de $\Delta_{\mathrm{inj}}/D$ sont égaux.} D'où :
\[
  \left\{
  \begin{array}{l}
    \text{géométries de drapeaux munies d'ordres}\\
    \text{totaux compatibles sur les simplexes}
  \end{array}
  \right\}
  \ \simeq\
  \left\{
  \begin{array}{l}
    \text{ensembles semi-simpliciaux}\\
    \text{satisfaisant (1)}
  \end{array}
  \right\} .
\]

\emph{Cellules duales} (page 19). Soit $D$ une géométrie de drapeaux munie
d'ordres compatibles, et $X = |D|$, avec ses simplexes fermés $X(d)$ et ses
simplexes ouverts $X^\circ(d)$, qui forment une partition de $X$. Pour un
sommet $F \in \Phi$, soit
\[
  |F| = \bigcup_{\sup I_d = F} X(d) , \qquad
  |F|^\circ = \bigcup_{\sup I_d = F} X^\circ(d) ,
\]
la réunion des simplexes dont $F$ est le \emph{dernier} sommet. Quand $D$ est
le complexe des chaînes d'un complexe cellulaire, ordonné par la dimension,
$|F|$ est la cellule fermée de $F$ subdivisée, et $|F|^\circ$ la cellule
ouverte. La page demande quand les $|F|$ sont les cellules fermées d'une
stratification de $X$ dont les $|F|^\circ$ sont les cellules ouvertes, et
établit :
\begin{enumerate}
  \item[a)] les $|F|^\circ$ forment une partition de $X$, puisque chaque
  simplexe ouvert a un seul dernier sommet, et chaque $|F|^\circ$ contient le
  sommet $F$ ;
  \item[b)] $|F|$ est l'adhérence de $|F|^\circ$, puisque $X(d)$ est celle de
  $X^\circ(d)$.
\end{enumerate}
Il reste la question qu'elle pose en marge, $|F| \smallsetminus |F|^\circ =
\bigcup_{F' < F} |F'|$ ? L'inclusion $\subset$ est vraie : un point de $|F|$
hors de $|F|^\circ$ est dans un simplexe ouvert dont le dernier sommet précède
$F$. L'égalité demande que tout simplexe de dernier sommet $F' < F$ soit face
d'un simplexe de dernier sommet $F$, ce qui n'est pas
automatique.\footnote{Cette discussion de l'égalité est la nôtre ; la page la
marque d'un point d'interrogation. Ce que sont les $|F|$ --- des boules ou
non --- est la question b) de la page 29.}

\pagerange{20}{24}
\section*{Ensembles semi-simpliciaux et ordre sur les sommets (pages 20 à 24)}

\emph{Le point de départ} (pages 20 et 21). Soit $D_*$ un ensemble
semi-simplicial, $D_n = D_*(\Delta_n)$. La catégorie fibrée qu'il définit sur
$\Delta_{\mathrm{inj}}$, à fibres discrètes, est ordonnée si et seulement si
$D_*$ satisfait (1). L'ensemble ordonné $D_* = \coprod_m D_m$ a alors la
propriété
\begin{enumerate}
  \item[(a)] pour tout $d \in D_m$, l'idéal $D_*(d) = \{\delta \leqslant d\}$
  est isomorphe à $\mathfrak{P}^*(\Delta_m)$,
\end{enumerate}
c'est donc une géométrie de drapeaux, et sa réalisation est
\[
  |D_*| = \varinjlim_{d \in D_*} |d| = \varinjlim_{(\Delta_m, d) \in
  \Delta_{\mathrm{inj}}/D_*} |\Delta_m| ,
\]
la réalisation de l'ensemble semi-simplicial ; chaque simplexe $|d|$ s'y
plonge, et $d \mapsto |d|$ est injective et croissante.\footnote{La page 21
dit que $|D_*|$ est « aussi la réalisation topologique de l'ensemble
quasi-simplicial $D_*$ ». La réalisation de l'ensemble ordonné $D_*$, au sens
de la page 4, est la subdivision barycentrique de celle de l'ensemble
semi-simplicial : les deux sont homéomorphes, canoniquement, sans être
égales. Le plongement de chaque simplexe tient à (1).}

\emph{La question.} La donnée d'ordres totaux compatibles sur les $I_d$
équivaut à la structure semi-simpliciale. Quand ces ordres viennent-ils d'un
seul ordre sur $D_0$ ?

\emph{L'ordre des sommets} (page 22). Sur $D_0$, posons
\[
  x \prec y \iff \text{il existe } z \in D_1 \text{ avec } s(z) = x,\ t(z) = y .
\]
\begin{enumerate}
  \item[(2°)] Cette relation est transitive si
  \item[(2 bis)] l'application $D_2 \to D_1 \times_{D_0} D_1$, $u \mapsto
  (d_2u, d_0u)$, le produit fibré étant pris au-dessus de $t$ et de $s$, est
  surjective.
\end{enumerate}
En effet, si $x \prec y$ par $z$ et $y \prec w$ par $z'$, un $u$ tel que $d_2 u
= z$ et $d_0 u = z'$ donne l'arête $z'' = d_1 u$, de $x$ à $w$. La condition
(2 bis) dit que tout couple d'arêtes composables est bordé par un
$2$-simplexe.\footnote{C'est la condition de remplissage des cornets internes
$\Lambda^2_1$, en existence seulement --- la forme en dimension 2 de la
condition des quasi-catégories, ici pour un ensemble semi-simplicial. Le nom
est le nôtre. La page note $\partial_0, \partial_1$ à la page 21 puis
$\delta_0, \delta_1, \delta_2$ ; sur $D_1$, son $\delta_0$ est le premier
sommet et son $\delta_1$ le second, sur $D_2$ son $\delta_i$ est l'arête
opposée au sommet $i$ ; les identités $\delta_0\delta_1 = \delta_0\delta_2$
et $\delta_1\delta_1 = \delta_1\delta_0$ qu'elle écrit sont justes dans cette
convention mixte, et ne le seraient pas avec les faces usuelles. On écrit $s$
et $t$ pour éviter l'ambiguïté. Une note en diagonale, au crayon bleu, propose
d'exiger aussi d'une fonction dimension $\delta(x) < \delta(y)$ quand $x \prec
y$ ; elle est en grande partie illisible.} La relation $x \prec x$ est
impossible par (1), puisque les deux sommets d'une arête sont distincts ;
sous (2°), $x \preccurlyeq y \iff (x = y \text{ ou } x \prec y)$ est donc un
ordre sur $D_0$. Il induit sur chaque $I_d$ l'ordre total donné, et satisfait
\begin{enumerate}
  \item[(b)] si $x \prec y$, il existe $z \in D_*$ qui majore $x$ et $y$ ;
  \item[(c)] il induit un ordre total sur chaque $I_d$.
\end{enumerate}

\emph{La réciproque} (pages 23 et 24). Soit $D_*$ un ensemble ordonné
satisfaisant (a), et $\preccurlyeq$ un ordre sur $D_0$ satisfaisant (b) et
(c). L'ordre total induit sur $I_d$ est fonctoriel en $d$, donc définit un
foncteur fibrant à fibres discrètes $D_* \to \Delta_{\mathrm{inj}}$, c'est-à-dire
un ensemble semi-simplicial, qui satisfait (1) ; et l'ordre $\prec_*$ que la
page 22 lit sur ses arêtes est l'ordre de départ --- c'est ce que (b)
assure. Les deux constructions sont inverses l'une de l'autre, et donnent :

\textbf{Proposition.}
\[
  \begin{array}{c}
  \left\{
  \begin{array}{l}
    \text{ensembles semi-simpliciaux satisfaisant (1)}\\
    \text{et la transitivité (2°)}
  \end{array}
  \right\}
  \\[1ex] \Big\updownarrow \\[1ex]
  \left\{
  \begin{array}{l}
    \text{ensembles ordonnés satisfaisant (a), munis}\\
    \text{d'un ordre sur } D_0 \text{ satisfaisant (b) et (c)}
  \end{array}
  \right\}
  \end{array}
\]
\footnote{La page écrit, à gauche, « satisfaisant (1) et (2) », le second
chiffre d'une lecture incertaine. C'est (2°), la transitivité, qui convient,
et non (2 bis), plus forte : l'ensemble semi-simplicial associé à $(D_*,
\preccurlyeq)$ n'a pas de raison de remplir les cornets, comme le montre le
bord d'un triangle $\{0 < 1 < 2\}$, où $0 \prec 2$ est une arête mais aucun
$2$-simplexe ne borde $(01, 12)$. À la page 24, la page écrit « $x_0 =
\delta_0(z_1)$, $y_0 = \delta_0(z_1)$ » ; le second est $\delta_1$, comme à la
page 23.}

\textbf{Remarque} (page 24). Si $\omega$ est un ordre sur $D_0$ satisfaisant
(b) et (c), l'ordre opposé les satisfait aussi, et donne une autre structure
semi-simpliciale sur les mêmes $D_n$, distincte de la première dès que $D_1
\neq \emptyset$. Les deux réalisations sont néanmoins canoniquement
homéomorphes, à travers la réalisation de l'ensemble ordonné $D_*$, qui ne
dépend d'aucun choix.

\pagerange{25}{25}
\section*{Premier exemple : les drapeaux d'un ensemble ordonné (page 25)}

Soit $\mathfrak{X}$ un ensemble ordonné, et $D_*$ l'ensemble de ses chaînes
finies non vides --- ses drapeaux ---, ordonné par inclusion. Alors $D_0 =
\mathfrak{X}$, (a), (b), (c) sont satisfaits, et l'ensemble semi-simplicial
obtenu est le nerf de $\mathfrak{X}$ privé de ses dégénérescences ; ses
réalisations sont canoniquement homéomorphes à celle de $\mathfrak{X}$.

Ces exemples satisfont de plus
\begin{enumerate}
  \item[(d)] tout $d$ est la borne supérieure de $I_d$,
\end{enumerate}
c'est-à-dire que $D_*$ est l'ensemble des simplexes d'un complexe
simplicial --- d'une « multiplicité simpliciale », dit la page. Mais (a) à (d)
ne caractérisent pas les drapeaux d'un ensemble ordonné : ils caractérisent
les complexes simpliciaux munis d'un ordre sur les sommets, total sur chaque
simplexe, et tel que deux sommets comparables soient joints par une arête. Les
drapeaux de $\mathfrak{X}$ sont ceux où, de plus, toute partie totalement
ordonnée est un simplexe.\footnote{La page commence par « les $(D_*, \prec)$
obtenus sont ceux pour lesquels on a, en plus de (a), (b), (c), la condition
(d) », puis se reprend : « Donc il faut prendre les multiplicités simpliciales
avec une structure d'ordre sur l'ensemble des sommets, induisant un ordre
total sur chaque simplexe, et \ldots\ (sur les simplexes seulement\ldots ».
Le bord d'un triangle $\{0 < 1 < 2\}$ satisfait (a) à (d) et ne vient d'aucun
ensemble ordonné. Le crochet final affirme que (a) + (d) équivaut à (a') +
(d), où (a') dit que dans $D_*$ toute famille finie majorée a une borne
supérieure ; l'implication de gauche à droite est juste, la réciproque non :
trois éléments minimaux sous un même élément, et rien d'autre, satisfont (a')
et (d) sans satisfaire (a). Il faut ajouter que la borne supérieure d'une
partie $S$ de sommets n'a pas d'autre sommet que ceux de $S$.}

\pagerange{26}{29}
\section*{Second exemple : ensembles à opérateurs et drapeaux de type $H$ (pages 26 à 29)}

Soit $G$ un groupe et $\sigma_0, \ldots, \sigma_n \in G$ tels que $\sigma_i$
et $\sigma_j$ commutent pour $|i - j| \geqslant 2$. Pour $K \subset \Delta_n$,
soit $G_K$ le sous-groupe engendré par les $\sigma_i$, $i \in K$, et, pour
$H \subset \Delta_n$ non vide, $N_H = G_{\Delta_n \smallsetminus H}$. Soit $X$
un $G$-ensemble ; posons
\[
  D(H) = X / N_H \qquad (\text{drapeaux de type } H) .
\]
Si $H \subset H'$, $N_{H'} \subset N_H$ donne une restriction $D(H') \to D(H)$,
et $H \mapsto D(H)$ est un préfaisceau sur $\mathfrak{P}^*(\Delta_n)$. Le
modèle est l'ensemble $X$ des drapeaux complets d'une décomposition, avec
l'action du groupe cartographique : $X/N_H$ est l'ensemble de ses drapeaux de
type $H$, comme le voulait l'axiome 4 de la page 3.\footnote{C'est la
construction des géométries de classes de Tits, $G/N_H$ quand $X = G$, et des
« systèmes de chambres ». Le dossier 88 (pages 4 et 5) fait la même
construction pour le groupe de Weyl et ses sous-groupes paraboliques, et
obtient le complexe de Coxeter. La page 26 prend $H \in \mathfrak{P}(\Delta_n)$,
vide compris ; on se limite aux parties non vides, les seules qui donnent des
drapeaux.} L'ensemble ordonné $D_* = \coprod_H D(H)$, avec $d \leqslant d'$ si
$H \subset H'$ et $d$ est la restriction de $d'$, satisfait (a), avec $I_d
\simeq H$ pour $d \in D(H)$ ; ses éléments minimaux forment $D_0 =
\coprod_i X/N_{\{i\}}$.

\emph{L'ordre sur $D_0$.} Pour $x \in D(\{i\})$, $y \in D(\{j\})$, posons $x
\prec y$ si $i < j$ et si $x$ et $y$ ont un majorant commun. Cette relation
satisfait (b) et (c) ; elle est transitive dès que
\[
  D(\{i, j, k\}) \longrightarrow D(\{i, j\}) \times_{D(\{j\})} D(\{j, k\})
  \qquad (i < j < k)
\]
est surjective, et c'est le cas par le lemme suivant.

\textbf{Lemme} (pages 27 et 28). Soient $H, H' \subset \Delta_n$ tels que le
plus grand élément de $H$ soit dans $H'$, et que $H \cap H'$ soit un segment
initial de $H'$. Alors
\[
  D(H \cup H') \longrightarrow D(H) \times_{D(H \cap H')} D(H')
\]
est surjective.\footnote{La page énonce le lemme sous l'hypothèse « $H \cap H'$
segment final de $H$ et segment initial de $H'$ », plus forte, et note en
marge de la démonstration n'avoir utilisé que les deux propriétés qu'on
énonce. Les segments s'entendent dans l'ordre induit sur $H$ et $H'$.}

\emph{Démonstration.} Avec $K = \Delta_n \smallsetminus H$ et $K' = \Delta_n
\smallsetminus H'$, la flèche est $X/G_{K \cap K'} \to X/G_K \times_{X/G_{K
\cup K'}} X/G_{K'}$, et il suffit que $G_{K \cup K'} = G_{K'} \cdot G_K$ :
si $(G_Kx, G_{K'}y)$ est dans le produit fibré, $y = k'kx$ avec $k' \in
G_{K'}$, $k \in G_K$, et $z = kx$ a les deux images voulues. Soit $i$ le plus
grand élément de $H$ ; il est dans $H \cap H'$, donc hors de $K \cup K'$, et
\[
  K \cup K' = K'_0 \sqcup K_0, \qquad K'_0 = (K \cup K') \cap [0, i-1],
  \quad K_0 = (K \cup K') \cap [i+1, n] .
\]
On a $K_0 \subset K$, car $[i+1, n]$ ne rencontre pas $H$ ; et $K'_0 \subset
K'$, car un $j < i$ qui serait dans $H'$ serait dans le segment initial $H
\cap H'$. Les indices de $K'_0$ et de $K_0$ diffèrent d'au moins deux, donc
$G_{K'_0}$ et $G_{K_0}$ commutent, et $G_{K \cup K'} = G_{K'_0} G_{K_0}
\subset G_{K'} G_K$.

Ainsi, pour \emph{tout} $G$-ensemble $X$, $(D_*, \prec)$ satisfait (a), (b),
(c) et la transitivité, et définit un ensemble semi-simplicial satisfaisant
(1), avec $D_m = \coprod_{\mathrm{card}\, H = m+1} X/N_H$. La page n'exige pas
des $\sigma_i$ qu'ils soient des involutions.

\emph{Programme} (page 29). La page s'arrête sur deux tâches : a) trouver,
en termes de l'ensemble ordonné $(D_*, \prec)$, des conditions pour que (2
bis) soit satisfaite, et plus généralement pour la surjectivité de $D(\Delta_n)
\to D(H) \times_{D(H \cap H')} D(H')$ quand $\Delta_n = H \cup H'$ ; b)
décrire la décomposition cellulaire de $|D_*|$ associée à l'ordre sur $D_0$ ---
la question de la page 19 --- et reconstituer $(D_*, <, \prec)$ à partir de la
structure d'espace décomposé en cellules et de ses structures coniques. Le
reste de la page est blanc.

\pagerange{30}{33}
\section*{Préfaisceaux sur les parties finies de $\mathbf{N}$ (pages 30 à 33)}

Soit $D$ un préfaisceau d'ensembles sur l'ensemble ordonné
$\mathfrak{P}^*(\mathbf{N})$ --- une paracarte au sens de la page 13 --- et
\[
  \Theta(D) = \coprod_{A} D(A), \qquad \Theta_n(D) =
  \coprod_{\mathrm{card}\, A = n+1} D(A) ,
\]
ordonné par $d \leqslant d'$ si $A \subset A'$ et $d$ est la restriction de
$d'$ : c'est l'ensemble ordonné associé à la catégorie fibrée de $D$, avec une
application croissante $\Theta(D) \to \mathfrak{P}^*(\mathbf{N})$. On
s'intéresse surtout au cas où les $D(A)$ sont finis et nuls hors des parties
de $[0, N]$ : $\Theta(D)$ est fini, et sa réalisation $|D|$ est un espace à
triangulation finie où $\Theta(D)$ se réalise comme un ensemble de parties.

L'ordre de $\mathbf{N}$ donne aussitôt une structure semi-simpliciale : si
$A$ a $m+1$ éléments, $u_A : \Delta_m \xrightarrow{\sim} A$ est la bijection
croissante, et pour $\sigma : \Delta_n \to \Delta_m$ strictement croissante,
$\sigma^* : D(A) \to D(u_A\sigma(\Delta_n))$ est la restriction. Faisant
varier $A$, on obtient $\sigma^* : \Theta_m(D) \to \Theta_n(D)$, et les
$\Theta_n(D)$ forment un ensemble semi-simplicial. C'est le cas $I =
\mathbf{N}$ des pages 16 à 18, l'ordre de $\mathbf{N}$ fournissant les ordres
totaux compatibles sur les simplexes.

Tout ceci vaut pour un ensemble ordonné $N$ quelconque, en remplaçant
$\mathfrak{P}^*(\mathbf{N})$ par l'ensemble $\mathrm{Dr}(N)$ des drapeaux de
$N$ --- ses parties finies non vides totalement ordonnées ---, qui est
l'ensemble ordonné associé au foncteur $\varphi_N(\Delta_n) =
\{\text{applications strictement croissantes } \Delta_n \to N\}$, c'est-à-dire
au nerf non dégénéré de $N$. Un préfaisceau $D$ sur $\mathrm{Dr}(N)$ donne
alors des applications croissantes
\[
  \Theta(D) \longrightarrow \mathrm{Dr}(N) \longrightarrow \Delta_{\mathrm{inj}} ,
\]
la seconde au sens des catégories fibrées. La page 33 s'arrête là, aux deux
tiers vide.

\pagerange{35}{37}
\section*{Cartes cellulaires de dimension $n$ et groupe cartographique (pages 35 à 37)}

La page 35 ouvre une nouvelle suite, sous un titre encadré : « Cartes
cellulaires $n$-dim et groupes cartographiques ».

\emph{Filtrations admissibles.} Une filtration $\emptyset = X_{-1} \subset X_0
\subset \cdots \subset X_n = X$ est \emph{$n$-admissible} --- ou
\emph{régulière}, dit la page 36 --- si, pour $0 \leqslant i \leqslant n$ :
\begin{enumerate}
  \item[a)] le découpage $\widetilde{X}_i = \mathrm{Déc}(X_i, X_{i-1})$ est une
  somme topologique de boules fermées de dimension $i$, dont la partie
  au-dessus de $X_{i-1}$ est le bord $\partial\widetilde{X}_i$, et
  $p_i : \widetilde{X}_i \to X_i$ est fini --- de sorte que $X_i \smallsetminus
  X_{i-1}$ est une somme de boules ouvertes de dimension $i$ ;
  \item[b)] pour $i \geqslant 1$, la filtration de $\partial\widetilde{X}_i$
  image inverse de $X_0 \subset \cdots \subset X_{i-1}$ par $q_{i-1} :
  \partial\widetilde{X}_i \to X_{i-1}$ est $(i-1)$-admissible, et les
  applications $q_{i-1}^{-1}(X_j) \smallsetminus q_{i-1}^{-1}(X_{j-1}) \to X_j
  \smallsetminus X_{j-1}$ sont étales.
\end{enumerate}
Autrement dit, par récurrence : une filtration $0$-admissible est un espace
discret, une filtration $1$-admissible un graphe topologique dont les arêtes
sont des segments attachés par leurs extrémités, et une filtration
$n$-admissible s'obtient d'une filtration $(n-1)$-admissible en attachant à
$X_{n-1}$ une réunion disjointe de boules fermées de dimension $n$ par une
application finie $q : \partial\widetilde{X}_n \to X_{n-1}$ telle que l'image
inverse de la filtration soit $(n-1)$-admissible et que $q$ soit étale sur
chaque strate.\footnote{La page écrit, dans cette dernière forme, « une
filtr.\ $n$-admissible de $\partial\widetilde{X}_n$ » ; par b), c'est
$(n-1)$-admissible. Le cas $0$ porte « nécessaire et \emph{suffisant} si
déf », sans doute pour un espace fini. Un premier essai de la définition, en
colonne de droite, est barré.} Ce sont des complexes CW dont les applications
d'attachement sont des homéomorphismes locaux sur chaque strate ; elles ne
sont pas injectives en général --- une boucle est un graphe $1$-admissible
---, de sorte que ce ne sont pas des complexes CW réguliers.\footnote{On ne
connaît pas de nom établi pour cette classe ; les réalisations des cartes
généralisées de Lienhardt (voir la note de la page 7) en sont proches. Le
rapprochement est le nôtre.}

\emph{Morphismes.} Un morphisme $f : X \to Y$ d'espaces à $n$-filtration
régulière est une application finie telle que $f^{-1}(Y_i) = X_i$ et que $X_i
\smallsetminus X_{i-1} \to Y_i \smallsetminus Y_{i-1}$ soit étale --- donc un
revêtement fini, $f$ étant propre. On note $\mathcal{F}_n$ leur catégorie et
$\mathcal{F}_n^\circ$ la sous-catégorie pleine des objets dont l'espace est
une variété topologique de dimension $n$, sans bord.\footnote{La page dit
« ss-catégorie pleine $\mathcal{F}_n^\circ$ des espaces \ldots\ à
$n$-filtration régulière », le mot manquant étant lu, avec doute,
« vérifiés ». On le comprend comme « variétés » : c'est ce qu'exigent la
construction de $\sigma_n$ à la page 37 (« si $X_n$ est une $n$-variété ») et
le théorème 1. Une note de marge, en grande partie illisible, parle de
foncteurs entre les $\mathcal{F}_i$ et les $\mathcal{F}_i^\circ$.}

\emph{Les foncteurs $E_n$.} Par récurrence sur $n \geqslant 1$, on construit
deux foncteurs qui rendent commutatif le carré

\begin{tikzcd}
  \mathcal{F}_n^\circ \arrow[r, "E_n^\circ"] \arrow[d, hook] & G_n\text{-ens}
  \arrow[d, "\text{oubli}"] \\
  \mathcal{F}_n \arrow[r, "E_n"] & G_{n-1}\text{-ens}
\end{tikzcd}

\begin{enumerate}
  \item[(1)] $n = 1$. $E_1(X_1) = \partial\widetilde{X}_1$, l'ensemble des
  extrémités des segments --- les demi-arêtes. L'application
  $\partial\widetilde{X}_1 \to \pi_0(\widetilde{X}_1)$, vers l'ensemble des
  arêtes, a ses fibres de cardinal 2 ; sa \emph{permutation caractéristique}
  $\sigma_0$ échange les deux bouts d'une arête. Si $X_1$ est une variété de
  dimension 1, $\partial\widetilde{X}_1 \to X_0$ a aussi ses fibres de
  cardinal 2, d'où $\sigma_1$, qui échange les deux demi-arêtes en un sommet.
  \item[(2)] $n \geqslant 2$. $\partial\widetilde{X}_n$, muni de la filtration
  image inverse, est une somme de sphères de dimension $n-1$ filtrées, donc un
  objet de $\mathcal{F}_{n-1}^\circ$, et l'on pose
  \[
    E_n(X_n) = E_{n-1}^\circ(\partial\widetilde{X}_n) ,
  \]
  muni de $\sigma_0, \ldots, \sigma_{n-1}$, c'est-à-dire de l'action de
  $G_{n-1}$ ; c'est la composée de $\partial : \mathcal{F}_n \to
  \mathcal{F}_{n-1}^\circ$ et de $E_{n-1}^\circ$. Si $X_n$ est une variété, le
  morphisme de $G_{n-2}$-ensembles $E_{n-1}(\partial\widetilde{X}_n) \to
  E_{n-1}(X_{n-1})$ a ses fibres de cardinal 2 --- une cellule de dimension
  $n-1$ est dans exactement deux cellules de dimension $n$ ---, et sa
  permutation caractéristique est $\sigma_n$ ; elle commute à $\sigma_0,
  \ldots, \sigma_{n-2}$.\footnote{La page écrit « il commute à $\sigma_1,
  \ldots, \sigma_{n-2}$ » ; il commute à tous les $\sigma_j$, $j \leqslant n -
  2$, $\sigma_0$ compris, puisque l'application est un morphisme de
  $G_{n-2}$-ensembles, et c'est ce que demande la présentation de $G_n$. Un
  premier essai du cas $n = 1$, au bas de la page 36, est barré ; il passait
  par le revêtement d'orientation des arcs.}
\end{enumerate}
Ainsi $E_n^\circ(X)$ est l'ensemble des drapeaux complets $(x_0, \ldots, x_n)$
de la décomposition, avec les opérations des pages 3 et 7 ; $E_n(X)$, pour un
$X$ qui n'est pas une variété, garde les opérations $\sigma_0, \ldots,
\sigma_{n-1}$ qui ont un sens à l'intérieur d'une cellule.

\pagerange{38}{40}
\section*{Le programme : théorèmes 1 à 4 (pages 38 à 40)}

\textbf{Théorème 1} (énoncé). Soit $n \geqslant 1$. Le foncteur
$E_n^\circ$ passe à la catégorie isotopique déduite de $\mathcal{F}_n^\circ$
et induit un foncteur pleinement fidèle
\[
  \mathcal{F}_n^\circ \longrightarrow (G_n\text{-ens}) ,
\]
dont l'image essentielle est formée de $G_n$-ensembles finis où les $\sigma_i$
opèrent sans point fixe.\footnote{« Isotypique » se lit à la page 38,
« isotopique » à la page 39 ; c'est la catégorie dont les morphismes sont les
classes d'isotopie. La page affirme que l'image essentielle est formée
\emph{des} $G_n$-ensembles où les $\sigma_i$ opèrent sans point fixe. C'est
juste pour $n \leqslant 2$ --- toute carte généralisée sans bord de dimension 2
se réalise sur une surface fermée ---, faux à partir de $n = 3$ : une orbite
de $\langle \sigma_0, \sigma_1, \sigma_2 \rangle$ décrit le bord d'une cellule
de dimension 3, qui peut être un tore aussi bien qu'une sphère, et une orbite
de $\langle \sigma_1, \sigma_2, \sigma_3 \rangle$ le link d'un sommet, qui doit
lui aussi être une sphère. La page 42 le reconnaît : caractériser l'image
essentielle « revient au même » que reconnaître les sphères. Le théorème doit
donc se lire, pour $n \geqslant 3$, comme une inclusion, l'égalité étant
l'objet du théorème 7. La finitude des $G_n$-ensembles, écrite à la page 7,
est omise ici.}

\emph{Cascades.} À $X_n \in \mathcal{F}_n$ on associe une cascade d'ensembles
à opérateurs
\[
  R_n(X_n) \to R_{n-1}(X_n) \to \cdots \to R_1(X_n) \to R_0(X_n),
  \qquad R_i(X_n) = E_i(X_i) ,
\]
$R_i$ étant un $G_{i-1}$-ensemble et $R_0(X_n) = X_0$ un simple ensemble ; la
flèche $R_i \to R_{i-1}$ s'obtient en appliquant $E_{i-1}$ à
$\partial\widetilde{X}_i \to X_{i-1}$, et c'est un morphisme de
$G_{i-2}$-ensembles. Soit $(\mathrm{Casc})_n$ la catégorie de ces cascades :
ce sont exactement les $n$-complexes combinatoires de la page 7, la condition
d'invariance de la page 7 étant celle qu'on vérifie sur les cascades d'espaces
filtrés.\footnote{L'identification est la nôtre ; la page 38 ne rappelle pas
la page 7, et ne précise pas les conditions de la catégorie (lue
« (Casca) »).}

\textbf{Théorème 2} (énoncé). Le foncteur $\varphi_n : \mathcal{F}_n \to
(\mathrm{Casc})_n$, de la catégorie isotopique déduite de $\mathcal{F}_n$, est
pleinement fidèle.

Il faut aussi un foncteur $\rho_n : (G_n\text{-ens}) \to (\mathrm{Casc})_n$
qui prolonge $\varphi_n$ sur $\mathcal{F}_n^\circ$ :

\textbf{Théorème 3} (énoncé). Pour $X_n \in \mathcal{F}_n^\circ$, de
$G_n$-ensemble $E = E_n^\circ(X_n)$,
\[
  R_n(X_n) \simeq E \ (\text{en oubliant } \sigma_n), \quad
  R_{n-1}(X_n) \simeq E/\langle \sigma_n \rangle, \quad \ldots, \quad
  R_0(X_n) \simeq E/\langle \sigma_n, \ldots, \sigma_1 \rangle ,
\]
les flèches de transition étant les passages au quotient, et
$R_{n-k}(X_n)$ gardant l'action de $\sigma_0, \ldots, \sigma_{n-k-1}$, qui
commutent à $\sigma_{n-k+1}, \ldots, \sigma_n$. On pose donc $\rho_n(E) = (E
\to E/\langle\sigma_n\rangle \to \cdots)$.\footnote{C'est la formule des bords
$\delta_{i-1}$ de la page 8. Elle tient à ce que, dans une variété, les
drapeaux complets qui prolongent un drapeau $(x_0, \ldots, x_j)$ sont ceux du
link de $x_j$, une sphère de dimension $n-j-1$, et sont permutés
transitivement par $\langle \sigma_{j+1}, \ldots, \sigma_n \rangle$ --- la
connexité forte des polytopes abstraits. La page appelle ce foncteur
$\psi_n$, comme celui de la page 40 ; on le note $\rho_n$, lettre qu'elle
emploie elle-même à la page 40.}

\emph{Le multicône} (page 40). En sens inverse, on cherche des foncteurs ---
adjoints ? --- $(\mathrm{Casc})_n \to \mathcal{F}_n$ et $(G_n\text{-ens}) \to
\mathcal{F}_n^\circ$. On définit $\psi_n : (\mathrm{Casc})_n \to \mathcal{F}_n$
par récurrence. Pour $n = 0$, $\psi_0(E_0)$ est $E_0$ muni de la topologie
discrète. En général, une cascade $C_n = (E_n \to E_{n-1} \to \cdots \to E_0)$
contient la cascade tronquée $C_{n-1} = (E_{n-1} \to \cdots \to E_0)$, qui
définit $X_{n-1} = \psi_{n-1}(C_{n-1})$. Les cellules de dimension $n$ seront
les orbites, $\pi_0(\widetilde{X}_n) \simeq E_n/G_{n-1}$ ; leur bord est
\[
  \partial\widetilde{X}_n = \psi_{n-1}\bigl(\rho_{n-1}(E_n)\bigr) ,
\]
et l'application d'attachement $\partial\widetilde{X}_n \to X_{n-1}$ s'obtient
en appliquant $\psi_{n-1}$ au morphisme de cascades $\rho_{n-1}(E_n) \to
C_{n-1}$ --- qui existe précisément grâce à la condition d'invariance de la
page 7. Alors
\[
  \psi_n(C_n) = \text{multicône de } \psi_{n-1}\bigl(\rho_{n-1}(E_n) \to
  C_{n-1}\bigr) :
\]
on attache à $X_{n-1}$ un cône sur chaque composante de
$\partial\widetilde{X}_n$, comme à la page 12.\footnote{« Multicône » est une
lecture douteuse, écrite au-dessus d'un mot biffé, et la lettre $\rho$ aussi.
La page note en marge « préciser ce point » ; l'existence du morphisme
$\rho_{n-1}(E_n) \to C_{n-1}$ à partir de l'invariance de la page 7 est
notre précision. Un cône sur une composante de $\partial\widetilde{X}_n$
n'est une boule que si cette composante est une sphère : en général,
$\psi_n(C_n)$ est une carte \emph{conique} au sens de la page 9, et non
cellulaire. C'est la même difficulté que dans la note au théorème 1.}

\textbf{Théorème 4} (annoncé). Description de $\psi_n$, et le cas échéant
propriété d'adjonction.

\pagerange{41}{42}
\section*{Théorèmes 5 à 7, et d'autres points (pages 41 et 42)}

\textbf{Théorème 5} (annoncé). Subdivisions barycentriques de
$\psi_n(C_n)$, et leur description par des ensembles
pseudo-simpliciaux.\footnote{« Pseudo-simpliciaux » est une lecture
incertaine.}

\emph{Remarque.} Se donner un homéomorphisme $X_n \simeq \psi_n\varphi_n(X_n)$
revient à se donner, pour toute cellule $\widetilde{X}_i^\alpha$, un
homéomorphisme $\widetilde{X}_i^\alpha \simeq
\mathrm{Cône}(\partial\widetilde{X}_i^\alpha)$ induisant l'identité sur le
bord.\footnote{Pour une boule, un tel homéomorphisme existe et il est unique à
isotopie près relativement au bord : c'est le « truc d'Alexander » (1923), qui
explique qu'on travaille dans la catégorie isotopique. Le nom n'est pas sur
la page.} La subdivision barycentrique se fait au-dessus d'une telle
structure ; itérée deux fois, elle donne une triangulation, dont il faudrait
décrire le complexe simplicial --- un ensemble et des parties finies stables
par inclusion --- en termes de la cascade $C_n$ : ce serait un théorème 6.

\textbf{Théorème 7} (annoncé, sans énoncé). Caractériser l'image essentielle
de $\mathcal{F}_n^\circ$ dans les $G_n$-ensembles, et par là celle de
$\mathcal{F}_n$ dans les cascades, « en termes combinatoires ! Voire, ce qui
revient au même : étudier si un espace triangulé est homéomorphe à une
sphère --- cela semble lié : l'hypothèse de Poincaré ! ».\footnote{On sait
depuis S.~P.~Novikov (années 1950--1960) qu'aucun algorithme ne reconnaît la
sphère de dimension $n \geqslant 5$ parmi les complexes simpliciaux finis ; en
dimension 3 la reconnaissance est décidable (Rubinstein, 1992 ; Thompson,
1994). Une caractérisation combinatoire de l'image essentielle ne peut donc
pas, pour $n \geqslant 5$, être une condition qu'on vérifie mécaniquement.
Références de mémoire, non vérifiées ; la page ne cite rien.}

\emph{Autres points à examiner} : les relations avec les ensembles ordonnés,
« polyèdres combinatoires », par les drapeaux ; la dualité dans
$\mathcal{F}_n^\circ$ ; les filtrations admissibles d'un produit d'espaces ;
les fibrations.\footnote{Pour la dualité, on observera que $\sigma_i \mapsto
\sigma_{n-i}$ est un automorphisme de $G_n$ ; sur les cartes généralisées il
échange une décomposition et sa décomposition duale. La remarque est la
nôtre.}

\pagerange{45}{45}
\section*{Aparté : petits ensembles et géométries finies (page 45)}

Une page de listing, sans lien apparent avec ce qui l'entoure, dit qu'un
petit ensemble fini « est » une géométrie finie : que se donner un ensemble à
$k$ éléments revient à se donner une telle géométrie, les isomorphismes se
correspondant. On lit, avec ce qu'on sait aujourd'hui :
\begin{itemize}
  \item 2 éléments : un torseur sous $\mathbf{F}_2$, ou une droite affine sur
  $\mathbf{F}_2$ ;
  \item 3 éléments : une droite affine sur $\mathbf{F}_3$, les vecteurs non
  nuls d'un plan sur $\mathbf{F}_2$, une droite projective sur $\mathbf{F}_2$
  --- de groupes $\mathrm{Aff}(1, \mathbf{F}_3) \simeq \mathrm{GL}(2,
  \mathbf{F}_2) \simeq \mathfrak{S}_3$ ;
  \item 4 éléments : un plan affine sur $\mathbf{F}_2$, une droite projective
  sur $\mathbf{F}_3$ --- $\mathrm{Aff}(2, \mathbf{F}_2) \simeq
  \mathrm{PGL}(2, \mathbf{F}_3) \simeq \mathfrak{S}_4$ ; et, avec une
  orientation, une droite affine sur $\mathbf{F}_4$, $\mathrm{Aff}(1,
  \mathbf{F}_4) \simeq \mathfrak{A}_4$ ;
  \item 5 éléments orientés : une droite projective sur $\mathbf{F}_4$,
  $\mathrm{PGL}(2, \mathbf{F}_4) \simeq \mathfrak{A}_5$ ;
  \item 5 éléments, sans orientation : par $\mathfrak{S}_5 \simeq
  \mathrm{PGL}(2, \mathbf{F}_5)$, un ensemble $E$ à 5 éléments définit une
  droite projective sur $\mathbf{F}_5$, à 6 points, que la page rapproche
  de l'icosaèdre.
\end{itemize}
\footnote{« Orientation » d'un ensemble fini s'entend comme le choix d'une
des deux classes d'ordres totaux modulo les permutations paires ; sur un
ensemble à 4 (resp.\ 5) éléments, il y a exactement deux structures de droite
affine sur $\mathbf{F}_4$ (resp.\ projective sur $\mathbf{F}_4$), échangées
par une transposition, ce qui justifie l'énoncé. La ligne sur 5 éléments est
très serrée ; on y lit « droite projective sur $\mathbf{F}_5$, les 6 pts
\ldots\ icosaèdre \ldots\ à $E$ », et c'est l'isomorphisme exceptionnel
$\mathfrak{S}_5 \simeq \mathrm{PGL}(2, \mathbf{F}_5)$ --- celui qui fait de
$\mathfrak{S}_5$ un sous-groupe transitif de $\mathfrak{S}_6$ --- qui donne
un sens à la phrase ; les six points correspondent aux six grandes diagonales
de l'icosaèdre, sur lesquelles agit $\mathfrak{A}_5$. Interprétation de notre
part. Le dossier 89 identifie de même les ensembles à trois éléments orientés
et les plans sur $\mathbf{F}_3$.}

La seconde moitié de la page compte les points des groupes de Chevalley sur
un corps fini. Pour $G$ réductif déployé de rang $r$, à $N$ racines
positives, de degrés $d_\lambda$,
\[
  \mathrm{card}\, G(\mathbf{F}_q) = q^N (q-1)^r \,\mathrm{card}\,
  (G/B)(\mathbf{F}_q) = q^N (q-1)^r \prod_\lambda \frac{q^{d_\lambda} -
  1}{q - 1} ,
\]
et $\mathrm{card}\,(G/B)(\mathbf{F}_q) = \sum_w q^{\ell(w)} =
P_{G/B}(\sqrt{q})$, la somme portant sur les cellules de Bruhat, où
\[
  P_{G/B}(T) = \frac{\prod_\lambda (1 - T^{2d_\lambda})}{(1 - T^2)^r}
\]
est le polynôme de Poincaré de la variété des drapeaux, ce qu'exprime aussi
$P_{BG} \cdot P_{G/B} = P_{BT}$.\footnote{La page écrit $P_{G/B} = (1 -
T^2)^r / \prod_\lambda (1 - T^{2d_\lambda})$, la fraction renversée ; la
ligne suivante, « $(1 : P_{G/B}(\sqrt{q})) = (q-1)^r/\prod_\lambda
(q^{d_\lambda} - 1)$ », est juste, et le contrôle sur $\mathrm{SL}_2$, $1 +
T^2$, tranche. Les indices $BG$, $BT$ (espaces classifiants) sont d'une
lecture incertaine.} Pour $\mathrm{PGL}_3$ --- « $GP(2)$ », $N = 3$, $r = 2$,
$d_\lambda = 2, 3$ --- cela donne $q^3(q-1)^2(q+1)(q^2+q+1)$. Un tableau
dressé à la main donne les ordres de $\mathrm{GL}(1)$, $\mathrm{Aff}(1)$,
$\mathrm{SL}(2)$, $\mathrm{PGL}(2)$, $\mathrm{GL}(2)$, $\mathrm{Aff}(2)$ sur
$\mathbf{F}_q$ pour $q = 2, 3, 4, 5$ ; toutes les valeurs sont justes, et ce
sont celles qu'il faut pour vérifier les identifications de la première
moitié : $6 = \mathrm{card}\, \mathrm{GL}(2, \mathbf{F}_2)$, $24 =
\mathrm{card}\, \mathrm{Aff}(2, \mathbf{F}_2) = \mathrm{card}\, \mathrm{PGL}(2,
\mathbf{F}_3)$, $12$, $60$, $120$.

\pagerange{47}{47}
\section*{Sujets de DEA (page 47)}

\begin{enumerate}
  \item[1)] La catégorie isotopique des $1$-complexes topologiques compacts
  à groupe fini d'automorphismes, avec plus de morphismes que les seuls
  isomorphismes.
  \item[2)] La catégorie isotopique des cartes cellulaires --- le programme
  des pages 35 à 42.
  \item[3)] La convexité du point de vue de la géométrie projective : pour
  $A$ convexe dans un espace projectif réel, existe-t-il un hyperplan qui ne
  le rencontre pas ?\footnote{La réponse dépend de la définition de la
  convexité projective, que la page ne donne pas : pour un convexe
  « propre », contenu avec son adhérence dans une carte affine, c'est oui par
  définition ; une droite projective entière, qu'on peut tenir pour convexe,
  rencontre tout hyperplan.}
  \item[4)] La géométrie des cubes et les espaces homogènes de
  $\mathfrak{S}_4$ --- que la fin du dossier aborde : le groupe des rotations
  du cube est $\mathfrak{S}_4$ (page 61).
\end{enumerate}

\pagerange{49}{50}
\section*{Cube combinatoire : le cube d'un ensemble à involution (pages 49 et 50)}

La chemise de la page 43 porte « cube combinatoire », au-dessus de
« 2-cartes cellulaires » biffé.

\emph{La maquette.} Soit $(\mathbb{B}, \sigma)$ un ensemble fini muni d'une
involution sans point fixe, $I = \mathbb{B}/\sigma$, $n = \mathrm{card}\, I$.
Posons
\[
  E = \mathrm{Hom}_{\{\pm 1\}}(\mathbb{B}, \mathbf{R}) = \{\varphi :
  \mathbb{B} \to \mathbf{R} \mid \varphi(\sigma b) = -\varphi(b)\} , \qquad
  C = \{\varphi \in E \mid |\varphi(b)| \leqslant 1 \ \forall b\} .
\]
Le choix d'une section de $\mathbb{B} \to I$ identifie $E$ à $\mathbf{R}^I$ et
$C$ au cube $[-1, 1]^I$ ; mais $C$ est défini sans ce choix. Un élément $b$
de $\mathbb{B}$ définit la facette $\{\varphi(b) = 1\}$, et $\sigma b$ la
facette opposée.\footnote{La page écrit $\mathbb{B} \hookrightarrow C
\hookrightarrow E$ sans dire comment $\mathbb{B}$ s'envoie dans $C$. Une
lecture naturelle envoie $b$ sur le centre de sa facette, la fonction qui vaut
$1$ en $b$, $-1$ en $\sigma b$ et $0$ ailleurs ; c'est la nôtre. La marge
de la page 49 dit que les éléments de $\mathbb{B}$ correspondent aux facettes
de codimension 1, et que $\sigma F$ est l'unique facette de codimension 1 qui
ne rencontre pas $F$.}

\textbf{Proposition} (page 49). Pour $J \in \mathrm{Sp}(\mathbb{B})$, soit
\[
  F(J) = \{\varphi \in C \mid \varphi(b) = 1 \ \forall b \in J, \
  \varphi(b) \neq \pm 1 \ \forall b \notin J \cup \sigma J\} ,
\]
une face ouverte de $C$. Alors $J \mapsto F(J)$ est une bijection de
$\mathrm{Sp}(\mathbb{B})$ sur l'ensemble des faces non vides de $C$, qui
renverse l'ordre --- les faces étant ordonnées par l'inclusion de leurs
adhérences. Le treillis des faces de $C$ est donc l'opposé de
$\mathrm{Sp}(\mathbb{B})$, auquel on ajoute un plus petit élément, la face
vide. La section vide donne $C$ lui-même, les sections totales les
sommets.\footnote{La page note $\mathrm{Spd}(\mathbb{B}/I)$ les sections
partielles et $\mathrm{Fac}(C)$, ou $\mathrm{Comb}(C)$, l'ensemble ordonné des
faces, le « polyèdre combinatoire » de $C$ ; elle appelle « facettes » toutes
les faces. On garde « face » pour une face quelconque et « facette » pour une
face de codimension 1.}

\textbf{Proposition} (page 49). Soit $p : \mathbb{B} \to I$ une surjection et
$\mathcal{P}(p)$ l'ensemble de ses sections partielles, ordonné par
inclusion. Le foncteur $p \mapsto \mathcal{P}(p)$, du groupoïde des
surjections vers les ensembles ordonnés, est pleinement fidèle. En effet, on
retrouve $\mathbb{B}$ comme l'ensemble des éléments minimaux de
$\mathcal{P}(p) \smallsetminus \{\emptyset\}$, et la relation $p(b) = p(b')$
comme : $b = b'$, ou aucune section totale ne contient à la fois $b$ et $b'$.

\textbf{Corollaire} (page 50). Pour deux maquettes $(\mathbb{B}, \sigma)$,
$(\mathbb{B}', \sigma')$ de cubes $C$, $C'$, les applications
\[
  \mathrm{Isom}\bigl((\mathbb{B}, \sigma), (\mathbb{B}', \sigma')\bigr)
  \xrightarrow{\ \alpha\ } \mathrm{Isom}_{\mathrm{aff}}(C, C')
  \xrightarrow{\ \beta\ } \mathrm{Isom}_{\mathrm{ord}}\bigl(\mathrm{Fac}(C),
  \mathrm{Fac}(C')\bigr)
\]
sont des bijections. En particulier
\[
  \mathrm{Aut}(C) \simeq \mathrm{Aut}(\mathbb{B}, \sigma) = \{\pm 1\}^I
  \rtimes \mathfrak{S}_I ,
\]
le centralisateur de $\sigma$ dans le groupe des permutations de
$\mathbb{B}$, d'ordre $2^n n!$.\footnote{C'est le \emph{groupe hyperoctaédral},
groupe de Weyl de type $B_n$ ; le nom n'est pas sur la page.}

\emph{Démonstration.} $\beta\alpha$ est bijective par les deux propositions ;
$\beta$ est injective, une application affine étant déterminée par son effet
sur les sommets ; donc $\beta$, surjective, est bijective, et $\alpha$ aussi.
La page ajoute qu'on peut remplacer les isomorphismes affines par les
isomorphismes projectifs.\footnote{C'est vrai pour $n \geqslant 2$, les
sommets du cube contenant alors un repère projectif, faux pour $n = 1$ : les
homographies de la droite qui fixent les deux extrémités d'un segment forment
un groupe à un paramètre qui préserve le segment.}

\textbf{Corollaire.} Les faces d'un cube sont des cubes.

\pagerange{51}{53}
\section*{Drapeaux du cube, antipodie, nombre de faces (pages 51 à 53)}

\textbf{Proposition} (page 51). $\mathrm{Aut}(C)$ opère simplement
transitivement sur l'ensemble $\mathrm{Drap}(C)$ des drapeaux complets $F_0
\subset F_1 \subset \cdots \subset F_n = C$ de faces ; c'est ce qu'on
appellera un polyèdre \emph{régulier}.

En effet, une face de dimension $d$ est contenue dans exactement $n - d$
faces de dimension $d+1$, de sorte que
\[
  \mathrm{card}\, \mathrm{Drap}(C) = 2^n \cdot n \cdot (n-1) \cdots 1 = 2^n
  n! = \mathrm{card}\, \mathrm{Aut}(C) ,
\]
et il suffit que l'action soit libre, ce qui vaut pour tout polyèdre convexe :

\textbf{Lemme} (page 51). Un automorphisme affine $u$ d'un polyèdre convexe
$C$ qui fixe un drapeau complet est l'identité.

\emph{Démonstration.} Par récurrence sur la dimension. Si $u$ fixe $(F_0,
\ldots, F_{n-1}, C)$, il est l'identité sur la facette $F_{n-1}$ par
hypothèse de récurrence, donc sur l'hyperplan $H$ qu'elle engendre. Un
automorphisme affine d'ordre fini --- $\mathrm{Aut}(C)$ est fini --- qui fixe
$H$ point par point est l'identité ou une symétrie par rapport à $H$ ; une
symétrie échange les deux demi-espaces que borde $H$ et ne peut respecter
$C$, qui est d'un seul côté.

\textbf{Corollaire.} Pour un cube $C_0$ de dimension $n$ muni d'un drapeau
$D_0$ et tout cube $C$ de dimension $n$, $u \mapsto u(D_0)$ est une bijection
$\mathrm{Isom}(C_0, C) \simeq \mathrm{Drap}(C)$. Ainsi $\mathrm{Drap}(C)$ est
un torseur à droite sous $G_0 = \mathrm{Aut}(C_0)$, et un torseur à gauche
sous $\mathrm{Aut}(C)$, qui est le centralisateur de l'action de $G_0$ ; et
\[
  C \longmapsto \mathrm{Drap}(C) , \qquad \{n\text{-cubes}\}
  \xrightarrow{\ \sim\ } \{G_0\text{-torseurs à droite}\} ,
\]
est une équivalence de groupoïdes. On peut prendre pour $(C_0, D_0)$ le cube
standard, de maquette $[1, n] \times \{\pm 1\}$, avec le drapeau défini par
les sections partielles $[1, i] \times \{+1\}$, $0 \leqslant i \leqslant n$,
d'où $G_0 \simeq \{\pm 1\}^n \rtimes \mathfrak{S}_n$.\footnote{« Produit ½
direct » : il écrit ½ pour « semi- ». Le mot « commutant » de la page 52 est
d'une lecture incertaine ; le centralisateur de l'action à droite de $G_0$
dans les permutations de $\mathrm{Drap}(C)$ est bien $\mathrm{Aut}(C)$.}

\emph{L'antipodie} (page 53). La symétrie centrale $-\mathrm{id}$ du cube
est, du point de vue combinatoire, $J \mapsto \sigma J$ sur les sections
partielles, et, du point de vue de la maquette, l'automorphisme $\sigma$ de
$(\mathbb{B}, \sigma)$. Pour toute face propre $F$, $F$ et $\sigma F$ sont
disjointes ; pour une facette, $\sigma F$ est l'unique facette qui ne la
rencontre pas. Le centre de $\mathrm{Aut}(C)$ est $\{1, \sigma\}$, de sorte
que $\sigma$ est son unique élément central non trivial.\footnote{La page
appuie ce dernier point sur le fait que le centre de $\mathfrak{S}_I$ est
trivial pour $n \geqslant 3$ ; l'énoncé vaut en fait pour tout $n \geqslant
1$ --- pour $n = 2$, $\mathrm{Aut}(C)$ est le groupe diédral d'ordre 8, de
centre $\{\pm 1\}$. Une proposition encadrée et barrée commençait à
énoncer : pour deux faces fermées non vides $F$, $F'$, $F \cap F' \neq
\emptyset$ ou $F \cap \sigma F' \neq \emptyset$, en demandant si cela
caractérise $\sigma$.}

\emph{Nombre de faces.} Le cube a $c_d = \binom{n}{d} 2^{n-d}$ faces de
dimension $d$, et
\[
  \sum_{d=0}^n (-1)^d c_d = \sum_{d=0}^n (-1)^d \binom{n}{d} 2^{n-d} = (2 -
  1)^n = 1 ,
\]
formule d'Euler--Poincaré qui vaut, souligne la page, pour tout polyèdre
convexe de dimension $n$, la face $C$ comptée.

\pagerange{54}{56}
\section*{Orientations, polyèdres réguliers (pages 54 à 56)}

\emph{L'orientation du cube.} L'ensemble à deux éléments des orientations de
$C$ est, par définition provisoire, celui de l'espace affine qu'il engendre,
et
\[
  \omega_C = \omega_E \simeq \omega_I \wedge \bigwedge_{i \in I}
  \mathbb{B}_i ,
\]
produit contracté de $\{\pm 1\}$-torseurs : une orientation de $E \simeq
\mathbf{R}^I$ est la donnée d'une orientation de l'ensemble $I$ --- un ordre
total à permutation paire près --- et d'un signe dans chaque fibre
$\mathbb{B}_i$, modulo un changement pair. L'antipodie $\sigma$ agit sur
chaque fibre, donc sur $\omega_C$ par $(-1)^n$.

\textbf{Proposition.} L'antipodie conserve l'orientation si et seulement si
$n$ est pair. \textbf{Corollaire.} Pour $n$ impair, $\mathrm{Aut}(C) \simeq
\mathrm{Aut}^+(C) \times \{\pm 1\}$.

\emph{Orientation combinatoire.} Pour un polyèdre convexe $C$ engendrant $E$,
l'ensemble $\omega_E$ se décrit en termes du seul ensemble ordonné
$\mathrm{Fac}(C)$. Un drapeau complet détermine une orientation --- celle de
la base $p_1 - p_0, \ldots, p_n - p_0$, pour des points $p_i$ choisis à
l'intérieur relatif de $F_i$, qui ne dépend que du drapeau ---, et une
orientation de $E$ revient à la partie $D^+$ des drapeaux « directs ». La
solution combinatoire : parmi les partitions $\mathrm{Drap}(C) = D^+ \sqcup
D^-$, ce sont celles que chaque $\sigma_i$, $0 \leqslant i \leqslant n-1$,
échange ; il y en a exactement deux.\footnote{La page écrit « les $\sigma_i$
($1 \leqslant i \leqslant n$) » ; sur les drapeaux d'un polyèdre de
dimension $n$, dont le dernier membre $C$ est fixe, ce sont $\sigma_0, \ldots,
\sigma_{n-1}$. Qu'il y ait exactement deux telles partitions demande que le
graphe des drapeaux soit connexe et biparti ; c'est le cas pour un polyèdre
convexe, puisque l'orientation géométrique change sous chaque $\sigma_i$.
Le dossier 89 (pages 19 et 20) énonce de même qu'un polygone combinatoire
connexe a exactement deux orientations. La page remplace partout
« drapeaux » par « repères » ; c'est le même objet.}

\textbf{Proposition} (page 55). Pour un polyèdre convexe $C$, les
homomorphismes
\[
  \mathrm{Aut}(C) \xrightarrow{\ \alpha\ } \mathrm{Aut}_{\mathrm{ord}}
  (\mathrm{Fac}(C)) \xrightarrow{\ \beta\ } \mathrm{Aut}_{\mathrm{ens}}\,
  \mathrm{Drap}(C)
\]
sont injectifs, et $\mathrm{Aut}_{\mathrm{ord}}(\mathrm{Fac}(C))$ agit
librement sur $\mathrm{Drap}(C)$. Pour $\alpha$ : un automorphisme est connu
sur les points extrémaux. Pour la liberté, soit $u$ un automorphisme de
$\Phi = \mathrm{Fac}(C)$ fixant un drapeau complet ; par récurrence, $u$ est
l'identité sur $\Phi_{\leqslant F_{n-1}}$. S'il est l'identité sur
$\Phi_{\leqslant x}$, $x$ une facette, et si $y$ est une facette qui partage
avec $x$ une face de dimension $n - 2$, il fixe $y$ --- seule autre facette
contenant cette face --- et un drapeau de $\Phi_{\leqslant y}$, donc est
l'identité sur $\Phi_{\leqslant y}$. Il reste à voir que la relation
« partager une face de dimension $n-2$ » engendre la relation grossière sur
les facettes, ce que la page renvoie à un « argument topologique de
connexité ».\footnote{C'est la connexité du graphe du polytope dual
(théorème de Balinski, 1961, qui la donne même $n$-fois connexe) ; le nom
n'est pas sur la page.}

\emph{Polyèdres réguliers et homogènes.} $C$ est \emph{régulier} si
$\mathrm{Aut}(C)$ est transitif --- donc simplement transitif --- sur
$\mathrm{Drap}(C)$, et \emph{homogène} si, pour $0 \leqslant i \leqslant n -
1$, toute face de dimension $i+1$ contient le même nombre $d_i$ de faces de
dimension $i$. Alors $\mathrm{card}\, \mathrm{Drap}(C) = d_0 d_1 \cdots
d_{n-1}$.\footnote{La définition de la page est embrouillée --- « $\forall F
\in \mathrm{Fac}_{i+1}(C)$, le cardinal de l'ens.\ des $y \in
\mathrm{Fac}_i(C)$ qui majorent $x$ soit $d_i$ » --- et suivie de « (car $d_n
= 1$) ». On compte $d_i$ vers le bas, ce qui rend la formule du nombre de
drapeaux juste : on choisit une facette parmi $d_{n-1}$, une face de
codimension 2 dans celle-ci parmi $d_{n-2}$, et ainsi de suite. Pour le cube,
$d_i = 2(i+1)$ et le produit vaut $2^n n!$.} Évidemment régulier implique
homogène. Avec
\[
  N = \mathrm{card}\, \mathrm{Aut}(C), \quad N_c = \mathrm{card}\,
  \mathrm{Aut}(\mathrm{Fac}(C)), \quad R = \mathrm{card}\, \mathrm{Drap}(C) ,
\]
on a $N \mid N_c \mid R$ --- un sous-groupe, puis une action libre --- et $C$
est régulier si et seulement si $N = R$, auquel cas $N = N_c = R = d_0 \cdots
d_{n-1}$.

\emph{Orientation d'un polyèdre convexe} (page 56) : une orientation de
l'espace affine $E_C$ qu'il engendre. Par ce qui précède, l'ensemble
$\omega(C)$ ne dépend, à isomorphisme canonique près, que de
$\mathrm{Fac}(C)$, d'où des applications $\mathrm{Isom}_{\mathrm{ord}}(\mathrm{Fac}(C),
\mathrm{Fac}(C')) \to \mathrm{Isom}(\omega(C), \omega(C'))$ et un
homomorphisme $\mathrm{Aut}_{\mathrm{ord}}(\mathrm{Fac}(C)) \to \{\pm 1\}$.
Si $C$ a une antipodie --- un centre de symétrie ---, sa signature est
$(-1)^n$.

\pagerange{57}{60}
\section*{Les classes de conjugaison de $\mathrm{Aut}(C_n)$ (pages 57 à 60)}

\emph{Réduction.} Un automorphisme $u$ du cube est un automorphisme de la
maquette ; il induit $u_I$ sur $I$, qui décompose $I$ en orbites $I_\alpha$,
et $u_I$ est connu à conjugaison près par les cardinaux de ces orbites. Avec
$\mathbb{B}_\alpha$ l'image inverse de $I_\alpha$, la situation est la somme
des $(\mathbb{B}_\alpha, u|\mathbb{B}_\alpha)$ : géométriquement, $C =
\prod_\alpha C_\alpha$ et $u = \prod_\alpha u_\alpha$. On est ramené au cas où
$u_I$ est un cycle sur $I$, de longueur $n$. Il y a alors deux cas :
\begin{enumerate}
  \item[a)] $\mathbb{B}$ est réunion de deux orbites de $u$, chacune une
  section de $\mathbb{B} \to I$ ; si $\varepsilon$ est l'ensemble de ces deux
  orbites, $\mathbb{B} \simeq \varepsilon \times I$ et $u = \mathrm{id}_\varepsilon
  \times u_I$ ;
  \item[b)] le « cas tordu » : $\mathbb{B}$ est une seule orbite, de longueur
  $2n$, et $u^n = \sigma$. La situation est déterminée à isomorphisme près
  par $n$.
\end{enumerate}
Dans des coordonnées, a) est un cycle de signe total $+1$ et b) un cycle de
signe total $-1$.

\textbf{Proposition} (page 58). Les classes de conjugaison de $\mathrm{Aut}(C_n)
= \{\pm 1\}^n \rtimes \mathfrak{S}_n$ sont paramétrées par deux suites
d'entiers $(\alpha_i^+)$, $(\alpha_i^-)$, $\alpha_i^+$ (resp.\
$\alpha_i^-$) étant le nombre de cycles de longueur $i$ de $u_I$ au-dessus
desquels $u$ a deux cycles de longueur $i$ (resp.\ un cycle de longueur
$2i$), soumises à
\[
  \sum_i i\,(\alpha_i^+ + \alpha_i^-) = n .
\]
\footnote{C'est la description classique par couples de partitions (« type
cyclique signé »), qu'on attribue à A.~Young (1930) ; le nom n'est pas sur la
page. La page lit « au-dessus desquels il y a un cycle de long.\ $i$ (resp.\
$2i$) » ; dans le cas non tordu il y en a deux.}

\emph{Faces fixes.} Les faces non vides fixées par $u$ correspondent aux
sections partielles stables par $u$. Au-dessus d'un cycle tordu, une telle
section est vide, puisque l'orbite contient les deux points de chaque fibre ;
au-dessus d'un cycle non tordu, elle est vide ou l'une des deux orbites. Il y
a donc $3^{\sum_i \alpha_i^+}$ faces fixes, $C$ comprise, et, si $f_d$ est le
nombre de faces fixes de codimension $d$ --- celles des sections à $d$
éléments ---,
\[
  \sum_{d=0}^n f_d\, t^d = \prod_{i \geqslant 1} (1 + 2t^i)^{\alpha_i^+} .
\]
\footnote{La page écrit « stables par $I$ » pour « stables par $u$ », et
indexe la somme par $i$ au lieu de $d$.}

\emph{Le cas $n = 3$} (pages 59 et 60). La page énumère les classes, avec en
marge le signe du déterminant et le cardinal de la classe :

\begin{itemize}
  \item un 3-cycle, non tordu : rotation d'ordre 3 autour d'une grande diagonale ($+$, 8) ;
  \item un 3-cycle, tordu : la précédente composée avec l'antipodie, ordre 6 ($-$, 8) ;
  \item 1 + 2, non tordus : symétrie par rapport à un plan diagonal ($-$, 6) ;
  \item 1 tordu, 2 non tordu : demi-tour autour d'un axe passant par deux milieux d'arêtes ($+$, 6) ;
  \item 1 non tordu, 2 tordu : quart de tour autour d'un axe passant par deux centres de faces ($+$, 6) ;
  \item 1 + 2, tordus : le précédent composé avec l'antipodie, ordre 4 ($-$, 6) ;
  \item 1 + 1 + 1, $\alpha_1^+ = 3$ : identité ($+$, 1) ;
  \item 1 + 1 + 1, $\alpha_1^+ = 2$ : symétrie par rapport à un plan de coordonnées ($-$, 3) ;
  \item 1 + 1 + 1, $\alpha_1^+ = 1$ : demi-tour autour d'un axe de coordonnées ($+$, 3) ;
  \item 1 + 1 + 1, $\alpha_1^+ = 0$ : antipodie ($-$, 1).
\end{itemize}

« En tout il y a 10 classes ? » : oui. C'est le nombre de couples de
partitions de poids total 3, et les cardinaux font $8 + 8 + 4 \cdot 6 + 1 + 3
+ 3 + 1 = 48$. Les cinq classes de signe $+$, de cardinaux $1, 8, 3, 6, 6$,
sont celles de $\mathfrak{S}_4$.\footnote{Le décompte de contrôle et le
rapprochement avec $\mathfrak{S}_4$ sont les nôtres ; la page met un point
d'interrogation après « 10 classes ». Pour le 3-cycle, la page note que la
relation entre cas tordu et non tordu, par composition avec l'antipodie, vaut
« chaque fois que la dimension $n$ est impaire » : pour un cycle de longueur
impaire, changer tous les signes change le signe total. Dans le cas
« 1 tordu, 2 non tordu », la page écrit « symétrie autour, composée du
précédent avec l'antipodie » : c'est bien un demi-tour.}

\pagerange{60}{61}
\section*{Le cube de dimension 3 et ses quatre diagonales (pages 60 et 61)}

\textbf{Proposition} (page 60). Soit $n \geqslant 3$, et $u$ un automorphisme
de $C$ tel que $u x = \pm x$ pour tout sommet $x$. Alors $u = \pm
\mathrm{id}$.\footnote{L'hypothèse $n \geqslant 3$, ajoutée en interligne, est
nécessaire : pour le carré, l'échange des deux coordonnées fixe deux sommets
et envoie les deux autres sur leurs opposés.}

Donc $\mathrm{Aut}(C)/\{\pm 1\}$ s'injecte dans le groupe des permutations de
$\Phi_0(C)/\sigma$, l'ensemble des $2^{n-1}$ grandes diagonales, et, pour $n$
impair, $\mathrm{Aut}^+(C) \simeq \mathrm{Aut}(C)/\{\pm 1\}$ aussi.\footnote{La
page écrit, pour la seconde injection, $\mathrm{Aut}^+(C) \hookrightarrow
\mathrm{Aut}(\Phi_0(C))$, sans le quotient par $\sigma$ ; l'injection dans
les permutations des sommets est vraie mais triviale, et c'est celle dans les
permutations des diagonales que la proposition donne, et que les corollaires
utilisent.}

\textbf{Corollaire} (page 61). Pour un cube de dimension 3,
$\mathrm{Aut}^+(C) \xrightarrow{\ \sim\ } \mathrm{Aut}_{\mathrm{ens}}
(\Phi_0(C)/\sigma)$, groupe des permutations des quatre grandes diagonales :
les deux ont 24 éléments.

\textbf{Corollaire.} $C \mapsto \Phi_0(C)/\sigma$ est une équivalence de la
catégorie des cubes orientés de dimension 3, avec leurs isomorphismes
directs, vers celle des ensembles à 4 éléments, avec leurs bijections. C'est
le quatrième sujet de DEA de la page 47, et le pendant des identifications de
la page 45 : un ensemble à 4 éléments « est » aussi un cube orienté.

\pagerange{61}{62}
\section*{Le cube gauche (pages 61 et 62)}

\emph{L'espace projectif découpé.} Soit $C = C(\mathbb{B}, \sigma)$ dans $E$.
Le quotient du bord $\partial C$ par l'antipodie est l'espace projectif réel
de dimension $n - 1$,
\[
  \partial C/\sigma \simeq (E \smallsetminus \{0\})/\mathbf{R}^* = P(E) ,
\]
chaque demi-droite issue du centre coupant $\partial C$ en un
point.\footnote{La page note $E^*/\mathbf{R}^*$, $E^*$ désignant $E$ privé de
l'origine et non le dual. Elle met sur le $C$ de $\partial C$ un petit signe,
chapeau ou flèche, qu'on ne reprend pas.} Soit $p : \partial C \to P(E)$ la
projection et $\Phi^{**}$ l'ensemble des faces propres non vides de $C$. La
décomposition de $\partial C$ en faces ouvertes passe au quotient :
\[
  \widetilde{C} = \partial C/\sigma = \coprod_{F \in \Phi^{**}/\sigma} p(F) .
\]
Chaque face fermée $\overline{F}$ s'injecte dans $P(E)$, puisque $\overline{F}
\cap \sigma\overline{F} = \emptyset$ (page 53), et
\[
  p(\overline{F}) \cap p(\overline{F'}) = p(\overline{F} \cap \overline{F'})
  \cup p(\overline{F} \cap \sigma\overline{F'}) ,
\]
réunion de deux fermés disjoints. Elle est vide si et seulement si les deux
intersections le sont ; elle n'est pas nécessairement connexe, et l'est si
l'une des deux est vide.\footnote{La seconde ligne de la formule de la page
61 est très surchargée et se lit $[p(\overline{F}) \cap p(\overline{F'})] \cup
[p(\overline{F}) \cap p(\overline{F'})]$, deux crochets identiques, avec un
trait qui pourrait être un $\sigma$ ; on écrit la formule juste. Le mot
« sinon » est d'une lecture douteuse, et une troisième ligne est biffée. Les
deux morceaux sont disjoints parce qu'un point commun serait dans $\overline{F'}
\cap \sigma\overline{F'} = \emptyset$, et leurs images le sont parce que $p$
est injectif sur $\overline{F}$.} Enfin
\[
  p(\overline{F}) \subset p(\overline{F'}) \iff \overline{F} \subset
  \overline{F'} \ \text{ ou } \ \overline{F} \subset \sigma\overline{F'} .
\]
\footnote{La page 62 écrit $p(\overline{F}) \subset p(\overline{F'}) \iff
\overline{F} \subset \overline{F'}$, qui est faux pour deux faces données :
une arête du carré est dans l'image de l'arête opposée. C'est vrai pour les
classes modulo $\sigma$, qui sont les cellules de $\widetilde{C}$ : $\overline{F}$,
connexe, est dans la réunion disjointe de $\overline{F'}$ et de
$\sigma\overline{F'}$ si et seulement s'il est dans l'une d'elles.}

On obtient une décomposition de $P_{\mathbf{R}}(E)$ en cellules
topologiques dont les cellules fermées sont plongées, deux cellules fermées
se rencontrant, si elles se rencontrent, en une cellule ou en la réunion de
deux cellules disjointes --- ce n'est donc pas un complexe polyédral, mais
c'est un complexe CW régulier. L'ensemble ordonné de ses cellules est
$\Phi^{**}/\sigma$, opposé de $\Phi^*/\sigma$, où $\Phi^*$ est l'ensemble des
sections partielles non vides de $\mathbb{B}$. La page dessine le cas $n = 2$ :
le carré de côtés $a, b, c, d$ devient une droite projective faite de deux
arcs, $\{a, c\}$ et $\{b, d\}$, qui se rencontrent en deux points.\footnote{Pour
$n = 3$ c'est l'\emph{hémi-cube}, polytope abstrait non sphérique que
McMullen et Schulte notent $\{4,3\}/2$ ; le nom n'est pas sur la page.}

\emph{Automorphismes des cubes gauches} ($n \geqslant 3$). La page considère
\[
  \mathrm{Aut}(C)/\{1, \sigma\} \longrightarrow \mathrm{Aut}_{\mathrm{proj}}
  (\widetilde{C}) \longrightarrow \mathrm{Aut}_{\mathrm{ord}}
  \bigl(\Phi(\widetilde{C})\bigr)
\]
et note que la première flèche est injective, la seconde évidemment aussi,
et demande si la composée est surjective. Pour $n$ impair, on récupère
$(\mathbb{B}, \sigma)$, avec sa décomposition cellulaire, comme le revêtement
d'orientation canonique de $\widetilde{C}$ --- $P(E)$ est alors de dimension
paire, donc non orientable, et son revêtement d'orientation est la sphère
$\partial C$ ---, et tout automorphisme de $\Phi(\widetilde{C})$ est induit
par un automorphisme direct de $C$. « Cas $n$ pair : à examiner\ldots », et le
dossier s'arrête là.\footnote{La réponse est oui pour tout $n \geqslant 3$, et
le revêtement universel remplace le revêtement d'orientation : un
automorphisme de l'ensemble ordonné des cellules d'un complexe CW régulier se
réalise par un homéomorphisme cellulaire ; pour $n \geqslant 3$, $\partial C
\to \widetilde{C}$ est le revêtement universel, l'homéomorphisme s'y relève,
et le relevé est un automorphisme combinatoire de $\partial C$, donc de $C$,
donc affine par le corollaire de la page 50. Ce raisonnement est le nôtre, et
la page ne le contient pas. Pour $n$ impair, $\mathrm{Aut}(C)/\{1, \sigma\}
\simeq \mathrm{Aut}^+(C)$, d'où la formulation de la page ; la lettre
$\mathcal{H}$ qu'elle emploie pour l'objet récupéré n'est pas définie et
désigne vraisemblablement le bord du cube.}

\end{document}
